\documentclass[11pt,a4paper]{article}

\usepackage[a4paper,margin=16mm,headheight=15pt]{geometry}
\usepackage{fontspec}
\usepackage[italian]{babel}
\usepackage{microtype}
\usepackage[table]{xcolor}
\usepackage{tikz}
\usetikzlibrary{arrows.meta,positioning,calc,shapes.geometric}
\usepackage[most]{tcolorbox}
\usepackage{tabularx,booktabs,array,longtable,multicol}
\usepackage{graphicx}
\usepackage{enumitem}
\usepackage{titlesec}
\usepackage{fancyhdr}
\usepackage{hyperref}
\usepackage{ragged2e}
\usepackage{pifont}
\usepackage{setspace}
\usepackage{needspace}

% Corpo di tradizione Bookman; font OpenType e alternative senza file esterni.
\defaultfontfeatures{Ligatures=TeX}
\IfFontExistsTF{TeX Gyre Bonum}
  {\setmainfont{TeX Gyre Bonum}}
  {\IfFontExistsTF{TeX Gyre Pagella}{\setmainfont{TeX Gyre Pagella}}{\setmainfont{DejaVu Serif}}}
\IfFontExistsTF{TeX Gyre Heros}{\setsansfont{TeX Gyre Heros}}{\setsansfont{DejaVu Sans}}
\IfFontExistsTF{TeX Gyre Pagella}{\newfontfamily\displayfont{TeX Gyre Pagella}}{\newfontfamily\displayfont{DejaVu Serif}}
\IfFontExistsTF{TeX Gyre Schola}{\newfontfamily\boxheadingfont{TeX Gyre Schola}}{\newcommand{\boxheadingfont}{\displayfont}}

\definecolor{ink}{HTML}{171717}
\definecolor{ash}{HTML}{E7E3DC}
\definecolor{mist}{HTML}{F4F2ED}
\definecolor{blood}{HTML}{5E1018}
\definecolor{iron}{HTML}{555555}
\definecolor{pale}{HTML}{FAF9F6}
\color{ink}

\hypersetup{colorlinks=true,linkcolor=blood,urlcolor=blood,pdftitle={Fate: Notti di Yharnam - Guida del GM - Il Progetto},pdfauthor={Materiale amatoriale per uso privato}}

\pagestyle{fancy}
\fancyhf{}
\fancyhead[L]{\small\textsc{Fate: Notti di Yharnam}}
\fancyhead[R]{\small\textsc{Guida del GM -- Il Progetto}}
\fancyfoot[C]{\small\thepage}
\renewcommand{\headrulewidth}{0.4pt}
\renewcommand{\footrulewidth}{0pt}

% ================================================================
% STILE YHARNAM -- seconda proposta
% Spunti visivi dalle librerie DND fornite dall'utente (Bergeron,
% Liu e Meyer): gerarchie tipografiche, note aperte, bordi da documento.
% Definizioni autonome: nessuna dipendenza da lib/dnd* o immagini.
% ================================================================
\definecolor{vellum}{HTML}{F4EEDF}
\definecolor{brass}{HTML}{91754B}
\definecolor{slatepaper}{HTML}{E8EDEE}

% Separatore rastremato: segue la larghezza locale, anche nei riquadri.
\newcommand{\yhdivider}{%
  \par\nobreak\noindent
  \begin{tikzpicture}
    \fill[brass!85] (0,0) -- (0,.7mm) -- (\linewidth,.18mm) -- cycle;
  \end{tikzpicture}\par\nobreak}

% Piccole losanghe sui bordi verticali delle letture.
% Sui riquadri spezzati, ornare solo l'inizio e la fine effettivi.
\newcommand{\yhreadingtop}{%
  \foreach \corner in {north west,north east}{%
    \fill[blood] ([yshift=1mm]frame.\corner)
      -- ([xshift=.7mm]frame.\corner)
      -- ([yshift=-1mm]frame.\corner)
      -- ([xshift=-.7mm]frame.\corner) -- cycle;}}
\newcommand{\yhreadingbottom}{%
  \foreach \corner in {south west,south east}{%
    \fill[blood] ([yshift=1mm]frame.\corner)
      -- ([xshift=.7mm]frame.\corner)
      -- ([yshift=-1mm]frame.\corner)
      -- ([xshift=-.7mm]frame.\corner) -- cycle;}}

% Estremita' oblique delle schede-documento, distinte dalle letture.
\newcommand{\yhpapersetup}{%
  \draw[brass,line width=.55pt]
    ([xshift=1mm,yshift=-2mm]frame.north west)
    -- ([xshift=3mm]frame.north west)
    -- ([xshift=8mm]frame.north west);
  \draw[brass,line width=.55pt]
    ([xshift=-8mm]frame.north east)
    -- ([xshift=-3mm]frame.north east)
    -- ([xshift=-1mm,yshift=-2mm]frame.north east);}
\newcommand{\yhpaperbase}{%
  \draw[brass,line width=.55pt]
    ([xshift=1mm,yshift=2mm]frame.south west)
    -- ([xshift=3mm]frame.south west)
    -- ([xshift=8mm]frame.south west);
  \draw[brass,line width=.55pt]
    ([xshift=-8mm]frame.south east)
    -- ([xshift=-3mm]frame.south east)
    -- ([xshift=-1mm,yshift=2mm]frame.south east);}

\tcbset{
  yhbase/.style={
    enhanced,breakable,arc=0mm,boxrule=0pt,
    left=4mm,right=4mm,top=3mm,bottom=3mm,
    before skip=10pt,after skip=10pt,
    fonttitle=\boxheadingfont\bfseries\scshape,
    toptitle=2mm,bottomtitle=1.5mm,
    before upper={\RaggedRight}},
  yhquiet/.style={
    yhbase,frame hidden,colback=mist,colframe=iron,
    colbacktitle=mist,coltitle=ink,
    fontupper=\sffamily\small},
  yhblood/.style={
    yhbase,frame hidden,colback=vellum,colbacktitle=vellum,
    colframe=blood,coltitle=blood,
    borderline north={1pt}{0pt}{blood},
    borderline south={.55pt}{0pt}{brass}}
}

\titleformat{\section}{\displayfont\LARGE\scshape\color{blood}}{}{0pt}{}[\yhdivider]
\titleformat{\subsection}{\displayfont\Large\scshape\color{blood}}{}{0pt}{}[\vspace{-1mm}{\color{brass!75}\titlerule[.6pt]}]
\titleformat{\subsubsection}{\boxheadingfont\normalsize\bfseries\scshape\color{ink}}{}{0pt}{}
\titlespacing*{\section}{0pt}{15pt}{9pt}
\titlespacing*{\subsection}{0pt}{13pt}{7pt}
\titlespacing*{\subsubsection}{0pt}{10pt}{4pt}
\setlist[itemize]{leftmargin=5mm,itemsep=1.5pt,topsep=3pt}
\setlist[enumerate]{leftmargin=6mm,itemsep=1.5pt,topsep=3pt}
\setlength{\parindent}{0pt}
\setlength{\parskip}{5pt}

\newtcolorbox{rulebox}[1][]{yhquiet,#1}
\newtcolorbox{readaloud}[1][]{
  yhbase,frame hidden,colback=vellum,colbacktitle=vellum,
  colframe=blood,coltitle=blood,
  borderline west={.85pt}{0pt}{blood},
  borderline east={.85pt}{0pt}{blood},
  left=5mm,right=5mm,top=3mm,bottom=3.5mm,
  fonttitle=\boxheadingfont\small\scshape,
  fontupper=\normalfont\linespread{1.08}\selectfont,
  title={Da leggere ad alta voce},
  overlay unbroken={\yhreadingtop\yhreadingbottom},
  overlay first={\yhreadingtop},
  overlay middle={},
  overlay last={\yhreadingbottom},#1}
\newtcolorbox{bloodbox}[1][]{yhblood,#1}
\newtcolorbox{gmbox}[1][]{
  yhquiet,colback=slatepaper,colbacktitle=slatepaper,
  title={Dietro la scena -- GM},#1}
\newtcolorbox{cluebox}[1][]{
  yhquiet,colback=vellum!60!white,colbacktitle=vellum!60!white,
  borderline north={.6pt}{0pt}{brass},
  title={Informazioni e indizi},#1}
\newtcolorbox{fatebox}[1][]{
  yhquiet,colback=ash!50!white,colbacktitle=ash!50!white,
  borderline west={1pt}{0pt}{iron},
  title={Fate -- Aspetti, azioni e compel},#1}
\newtcolorbox{clockbox}[1][]{
  yhblood,colback=blood!4!white,colbacktitle=blood!4!white,
  fontupper=\sffamily\small,
  title={Evoluzione -- Se il tempo avanza},#1}
\newtcolorbox{overviewbox}{
  yhquiet,colback=pale,left=2mm,right=2mm,top=2mm,bottom=2mm,
  before upper={\RaggedRight\sffamily\small
    \renewcommand{\arraystretch}{1.15}
    \rowcolors{1}{pale}{slatepaper!65!white}}}
\newtcolorbox{routebox}{
  yhquiet,colback=vellum!45!white,
  borderline north={.55pt}{0pt}{brass},
  borderline south={.35pt}{0pt}{brass}}
\newtcolorbox{handoutbox}[2]{
  yhbase,frame hidden,colback=vellum!70!white,
  colbacktitle=vellum!70!white,coltitle=ink,
  borderline north={.6pt}{0pt}{ink!65},
  borderline south={.6pt}{0pt}{ink!65},
  left=5mm,right=5mm,toptitle=3mm,bottom=4mm,
  title={\colorbox{blood}{\strut\textcolor{white}{%
    \sffamily\small\bfseries\enspace #1\enspace}}\quad #2},
  fontupper=\sffamily\small,
  overlay unbroken={\yhpapersetup\yhpaperbase},
  overlay first={\yhpapersetup},
  overlay middle={},
  overlay last={\yhpaperbase}}

\newcommand{\nodehead}[3]{%
\clearpage
{\titleformat{\section}{\displayfont\huge\scshape\color{blood}}{}{0pt}{}[\yhdivider]%
\section{#1}}
\begin{bloodbox}[title={Il nodo al tavolo}]
\textbf{Ruolo:} #2\par\smallskip
\textbf{Domanda del nodo:} #3
\end{bloodbox}}
\newcommand{\aspect}[1]{\textit{\textcolor{blood}{#1}}}
\newcommand{\tickbox}{\ding{113}}
\newcommand{\linkto}[1]{\begin{routebox}\textbf{Collega a:} #1\end{routebox}}
\fancyfoot[C]{\small\color{brass}\(\diamond\)\quad
  \textcolor{ink}{\thepage}\quad\(\diamond\)}

\begin{document}

\begin{titlepage}
\thispagestyle{empty}
\begin{tikzpicture}[remember picture,overlay]
  \draw[line width=1.2pt,blood] ([xshift=9mm,yshift=-9mm]current page.north west) rectangle ([xshift=-9mm,yshift=9mm]current page.south east);
  \draw[line width=0.35pt,iron] ([xshift=12mm,yshift=-12mm]current page.north west) rectangle ([xshift=-12mm,yshift=12mm]current page.south east);
  \node[opacity=.07] at (current page.center) {\fontsize{210}{210}\selectfont\textcolor{blood}{\ding{69}}};
\end{tikzpicture}
\vspace*{31mm}
\begin{center}
{\displayfont\fontsize{32}{37}\selectfont\bfseries FATE}\par
\vspace{2mm}
{\displayfont\fontsize{23}{28}\selectfont\bfseries NOTTI DI YHARNAM}\par
\vspace{6mm}
{\displayfont\fontsize{18}{22}\selectfont\color{blood}\bfseries IL PROGETTO}\par
\vspace{13mm}
{\large Guida del Game Master}\par
\vspace{20mm}
\begin{minipage}{.72\textwidth}
\centering\itshape
Una mini-campagna investigativa a nodi per 2--3 sessioni. I personaggi tornano a Yharnam per cercare un uomo scomparso e scoprono una verità che nessuna fazione dovrebbe possedere interamente.
\end{minipage}
\vfill
{\small Materiale amatoriale non ufficiale per uso privato. La continuità dell'ambientazione è volutamente elastica.}
\end{center}
\end{titlepage}

\tableofcontents
\clearpage

\section{1. La campagna in breve}

\begin{bloodbox}
\textbf{Obiettivo apparente:} ritrovare Elias Ward, scomparso da circa una settimana.\\
\textbf{Obiettivo emergente:} ricostruire l'indagine che Elias stava conducendo.\\
\textbf{Verità finale:} una cerchia interna alla Chiesa usa un reperto organico pthumeriano per trattare il sangue dei malati. I benefici temporanei spingono Celestine a imporre prove sempre più intense, usando il regolatore di Elias e cercando un contatto con il presunto Grande Essere d'origine.\\
\textbf{Scelta finale:} decidere il destino di Elias, dell'organismo e delle conoscenze raccolte.
\end{bloodbox}

\subsection{Tono e scala}
Questa non è una storia su come salvare Yharnam. La degenerazione provocata dal Sangue Antico è già troppo vasta per essere arrestata dai protagonisti. La vittoria possibile è salvare Elias, sventare il pericolo esperimento di Medre Celestine "Il Progetto".

Le apparizioni delle bestie devono servire a mettere pressione. Gli scontri dovrebbero essere rari, rapidi e collegati sempre a qualcosa.

\subsection{Struttura suggerita}
\begin{tabularx}{\linewidth}{>{\bfseries}p{25mm} X X}
Parte 1 & Dov'è Elias? & I giocatori capiscono che non è fuggito casualmente ma stava indagando, seguono le tracce ed acquisiscono le prime informazioni.\\
Parte 2 & Chi sta muovendo cosa? I giocatori acquisiscono conoscenza delle macchinazioni in moto, delle intenzioni di alcuni NPC e delle fazioni\\
Parte 3 & Ricollegare i punti e sventare il Progetto & I giocatori ora devono trovare Elias e decidere cosa fare per provare a salvare capra e cavoli\\
\end{tabularx}

\section{2. Principi di conduzione}


\subsection{Gli indizi fondamentali non richiedono un tiro}
Quando i PG cercano nel posto giusto o fanno una domanda sensata, ottengono l'informazione necessaria. I tiri stabiliscono rapidità, discrezione, precisione, vantaggi creati e complicazioni.

\subsection{Indizi ridondanti}
Non affidare un passaggio obbligato a un solo documento o personaggio. Se una pista viene chiusa, rendi disponibile una fonte alternativa già motivata.

\subsection{Ogni PNG conosce tre livelli diversi}
\begin{itemize}
\item \textbf{Sa}: fatti osservati direttamente.
\item \textbf{Crede}: interpretazione, spesso incompleta o errata.
\item \textbf{Nasconde}: informazione che non offre senza pressione, fiducia o scambio.
\end{itemize}

\subsection{Le fazioni sono alleati instabili}
La Chiesa, i Cacciatori e l'Accademia possono aprire porte e risolvere problemi. Sono ottimi spunti per il background dei PG, non sono ne buoni ne cattivi in senso assoluto, ognuna ha le sue verità ed i propri obiettivi. I PG possono chiedere ed ottenere aiuto, aumentando l'influenza della fazione coinvolta nella trama.
\section{3. La verità del master}

\begin{bloodbox}
Questa sezione fissa i fatti della campagna. Le convinzioni dei personaggi sono indicate separatamente.
\end{bloodbox}

\subsection{Celestine: risultati, convinzioni e deriva}
Madre Celestine Vey ha diretto per anni un reparto dell'Ospedale Vecchio. Dopo il ritrovamento di un reperto organico nei sotterranei pthumeriani, ipotizza che possa avere proprietà terapeutiche come il Sangue Antico. Lo conserva vivo in una vasca e costruisce una macchina che fa circolare il sangue dei pazienti attraverso il sistema contenente il reperto, per poi restituirlo al paziente: una sorta di dialisi.

Le prime prove sono volontarie, con rischi e limiti dichiarati, sedute brevi e rispetto del rifiuto dei malati. Alcuni pazienti tornano per poco tempo a riconoscere i propri cari e a parlare con lucidità, poi ricadono nello stato precedente.

Celestine ritiene che il reperto appartenesse a un Grande Essere e che il beneficio dipenda da un legame con esso. Seguendo testi pthumeriani, fa incidere una runa sulla vasca per favorire quel contatto. Quando il beneficio svanisce, conclude che il legame sia troppo debole. Aumenta durata, frequenza e portata dei trattamenti. Progressivamente minimizza i rischi, occulta ricadute ed effetti avversi, ignora le richieste di fermarsi e impone nuove sedute. Il regolatore di Elias viene aggiunto a una macchina già usata nelle prime prove.

\textbf{Fatti accertati:} il trattamento produce miglioramenti temporanei; i malati tornano poi nello stato precedente. Le prove attuali sono coercitive e i rapporti vengono manipolati.\par
\textbf{Convinzione di Celestine:} più sangue trattato e più sedute significherebbero un contatto più intenso e, infine, una cura stabile. Ogni beneficio conferma la teoria ai suoi occhi; ogni ricaduta giustifica un ulteriore aumento.\par
\textbf{Frattura della chiesa:} La maggioranza della Chiesa conosce un programma emoterapico con risultati promettenti, non la deriva delle prove o l'obiettivo ultraterreno. Celestine potrebbe trovare nemici potenti nella sua stessa organizzazione se venisse a galla quello che sta facendo.



\subsection{I tre elementi del Progetto: il reperto organico, la macchina ed il regolatore}
\begin{description}[leftmargin=5mm,style=nextline]
\item[Reperto organico.]
È il materiale vivente ritrovato nei sotterranei pthumeriani, conservato nella vasca. La sua appartenenza a un Grande Essere è l'ipotesi di Celestine. Potrebbe davvero aprire terribili canali di comunicazione con i grandi esseri in ascolto.
\item[Macchina.]
Comprende vasca, pompe e condotte. Durante una seduta il sangue esce dal paziente, attraversa il sistema con il reperto e ritorna al paziente. Mantiene anche le condizioni necessarie alla conservazione del reperto.
\item[Regolatore.]
È l'invenzione di Elias, originariamente destinata alle trasfusioni. Silas ha venduto i disegni, un'officina ecclesiastica ha costruito il pezzo. Aggiunto alla macchina dopo le prime prove, stabilizza la pressione e consente di aumentare la portata, facendo scorrere il sangue più velocemente. Celestine lo considera un mezzo per intensificare il contatto.
\end{description}

\textbf{Il ruolo dei pazienti:} i sei pazienti  trasferiti nell'ala inferiore vengono sottoposti a sedute sempre più insistenti. Il ritorno dei sintomi alimenta l'escalation di Celestine.

\begin{gmbox}
Il ruolo del sovrannaturale è lasciato nelle mani del master, può essere solo uno spauracchio o essere usato in modo più intenso, come modo per far finire la campagna se i PG decidono di far continuare l'esperimento contattando così davvero un grande essere, ad esempio. 
\end{gmbox}

\subsection{Simboli del manuale}
\textbf{Marchio amministrativo:} un cerchio con una freccia verso il basso identifica le commesse riservate gestite dalla cerchia di Celestine. Collega compratore dei disegni, forniture e autorizzazioni.

\textbf{Runa sulla vasca:} quattro bracci uncinati intorno a un centro vuoto. Nei testi consultati da Miriam è associata ad ascolto inteso come percezione, crescita e possibilità di contatto. Il disegno allegato alla commessa delle vetriature 44-12 (H06b) ne richiede l'incisione sulla vasca. È un tentativo di Celestine, il cui funzionamento non è dimostrato.

%\textbf{Percorso dell'indizio:} il giorno 18 Elias trova l'allegato in Archivio e copia soltanto la runa. Silas lo indirizza a Miriam, che il giorno 20 vede quella copia isolata: non il disegno della vasca, le commesse o i pazienti. H11 riporta una lettura dei testi. I PG possono trovare l'allegato H06b e mostrare a Miriam altre prove; aggiorna le proprie ipotesi soltanto in base a ciò che riceve.

\subsection{I progetti di Elias e la scelta di Silas}
Elias inventò un regolatore per trasfusioni prolungate. Due anni fa scoprì che era stato incorporato in strumenti usati per torturare prigionieri. Bruciò lo studio con i relativi disegni, poi divenne operaio alla Fonderia. La Chiesa sequestrò i frammenti dei prototipi, insufficienti però da soli a ricostruire il regolatore.

Dopo il rifiuto di Elias alla richiesta di un intermediario della Chiesa di realizzare il regolatore per uso medico, la Chiesa cercò i disegni altrove. Sette mesi fa Silas vendette una copia di lavoro che aveva conservato. Aveva debiti ed il compratore ecclesiastico prometteva un impiego ospedaliero: Silas si convinse che quella vendita potesse essere utile. La copia completò i disegni sequestrati e permise di costruire il regolatore con cui potenziare la macchina già in uso. Elias lo scopre durante l'indagine. Silas teme di essere considerato responsabile delle vittime e, prima di confessare, ammette solo vecchie commesse per la Chiesa.

\subsection{Topografia ed accesso al Progetto}
Il Progetto occupa l'ala inferiore riservata nei sotterranei dell'Ospedale Vecchio. L'ala comprende le stanze dei pazienti, un corridoio di servizio e la sala del Progetto con macchina e vasca.

Il Laboratorio Alchemico è un edificio adiacente all'Ospedale Vecchio e svolge normali preparazioni alchemiche per l'Ospedale. Questa attività reale giustifica personale e consegne e offre copertura alle forniture del Progetto. Il deposito interno riceve sia materiali per il lavoro ordinario sia carichi destinati all'ala inferiore riservata. Il carro dalla Fonderia consegna al deposito e da quel settore l'ascensore merci scende fino ad una breve galleria che raggiunge il corridoio dell'ala inferiore.

\begin{description}[leftmargin=5mm,style=nextline]
\item[Scala dell'Ospedale]
È l'accesso ordinario per personale e pazienti assistiti. Una porta chiusa e un sorvegliante ne controllano l'uso: serve un'autorizzazione. La scala conduce alle stanze dei pazienti, collegate alla sala del Progetto dal corridoio. Oren e Maud ne conoscono la posizione, ma non possono autorizzare l'ingresso di estranei.
\item[Ascensore del deposito.]
Si trova nel settore riservato del deposito interno al Laboratorio Alchemico, dietro una porta controllata. È noto agli addetti, non accessibile ai visitatori. La chiave è custodita dall'addetto all'ascensore. La chiave di Silas apre soltanto l'ingresso di servizio del Laboratorio che dà sul deposito, non l'ascensore. La chiave abilita il funzionamento e non deve essere tenuta da chi viaggia. Elias non la possiede: il giorno 26 usa l'ascensore già abilitato per un trasferimento.
\item[Scarico delle fogne.]
Lo scarico conduce dal vano di manutenzione adiacente alla sala del Progetto alle fogne cittadine, l'acqua defluisce verso le fogne. Una grata chiusa e fissata in posizione lascia passare l'acqua e deve essere superata per entrare. Quando non usato dal lavaggio il condotto è percorribile con difficoltà in senso opposto al deflusso, non è un accesso ordinario del personale. Alla quarta campana notturna il lavaggio lo rende impraticabile per circa un'ora. H14 indica il collegamento, la grata, la direzione e l'orario. Resta al GM definire come i PG possano percorrewre le fogne e raggiungere la grata.
\end{description}

\subsection{Forniture e falsificazioni}
Un'officina ecclesiastica prepara i componenti; la Fonderia esegue le finiture di precisione e li rispedisce. Il deposito interno al Laboratorio Alchemico smista anche le forniture del Progetto: l'ascensore merci le porta all'ala inferiore. Le consegne al Laboratorio sono ufficiali, l'attività alchemica ordinaria copre la destinazione finale e l'impiego di quei carichi. L'Ospedale fornisce personale, materiali sanitari e pazienti attraverso l'ala inferiore. L'Archivio conserva ordini e coperture.

\textbf{La cassa 17 e la 17-B sono la stessa cassa.} Il contabile della Chiesa ha descritto un regolatore medico come un ricambio industriale ordinario, copiando il peso standard di 412 libbre. Il contenuto reale pesa 391 libbre e un quarto; la bilancia grande arrotonda a 391. Per non correggere pubblicamente la falsa descrizione, il contabile chiude la voce 17 come sospesa e fa uscire il collo con il codice riservato 17-B. H04 è il foglio di pesatura originale conservato da Marta; H05 è la ricevuta di uscita del vetturale T. Harrow. Il peso e l'annotazione di ricodifica permettono il confronto. Questa è una prova che dimostra un occultamento amministrativo.

\subsection{Malattia e conoscenze di Elias}
Quattro mesi fa Elias ha ricevuto trasfusioni dopo un'ustione alla mano destra. Ha continuato il trattamento per tornare al lavoro. I primi sintomi compaiono sedici giorni prima dell'arrivo dei PG; nove giorni prima Oren li registra in H08. La malattia precede l'ingresso nei sotterranei.

Prima di entrare nella sala del Progetto Elias conosce l'anomalia dei pesi, la vendita dei disegni e la consegna al Laboratorio Alchemico: sospetta l'uso non autorizzato del regolatore. Lo riconosce direttamente soltanto il giorno 26, nella macchina usata per trattare il sangue dei pazienti. Elias sa di essere malato da prima dell'ingresso, ma non sa quanto tempo gli resti.

\subsection{Cronologia causale}
\textbf{Riferimento dei documenti:} oggi è il giorno 28 del mese corrente. Le date 12--15 in H04 indicano i giorni del mese. Le ore indicate dalle campane cittadine scandiscono i turni di lavoro.

\begin{longtable}{>{\bfseries}p{29mm} p{126mm}}
Due anni fa & Elias brucia lo studio dopo aver scoperto che i suoi strumenti vengono impiegati sui prigionieri. La Chiesa si presenta li e sequestra materiale incompleto. Silas conserva una copia dei disegni del regolatore.\\
Prima della vendita & Celestine ottiene il reperto e avvia prove volontarie a bassa portata. Osserva miglioramenti temporanei e ricadute; fa incidere la runa e decide di potenziare il trattamento. Un intermediario della Chiesa chiede a Elias di costruire il regolatore per uso medico. Elias rifiuta per timore di nuovi abusi.\\
Sette mesi fa & Silas vende la copia dei disegni al compratore ecclesiastico per saldare debiti, accettando una giustificazione di impiego in ambito medico. Il compratore usa il marchio amministrativo che Silas potrà riconoscere.\\
Sei mesi fa & Celestine amplia l'impianto e prolunga le prove, in attesa del regolatore. Il rispetto del consenso si deteriora: i rifiuti dei pazienti sono ignorati, i rischi minimizzati ed i  referti vengono lasciati incompleti. Le forniture passano per Fonderia e Laboratorio, l'Archivio conserva gli ordini.\\
Quattro mesi fa & Elias riceve sangue per l'ustione alla mano. Eleanor sa di queste cure; Oren, in ospedale, ne conserva la registrazione.\\
Giorno 12 & Compaiono insonnia e irritabilità. Elias le attribuisce al lavoro e continua a lavorare.\\
Giorno 13 & Una spedizione dello stesso committente indicato dall'intermediario attira l'attenzione di Elias. Pesa la cassa 17 con due strumenti e rileva la discrepanza. Scrive H01 senza aver visto il contenuto; sospetta che la commessa sia proseguita senza di lui. La cassa resta alla Fonderia fino alla partenza del giorno 16.\\
Giorno 14 & Elias chiede a Silas se siano rimaste copie dei vecchi disegni. Silas confessa di averne venduta una alla Chiesa. Questo rafforza il sospetto di Elias, ma non dimostra che la cassa contenga il suo regolatore.\\
Giorno 16 & La cassa 17 parte come collo 17-B sul carro di T. Harrow. Marta registra Elias come presente mentre lui si allontana. Elias segue il carro per verificare la destinazione dichiarata e vede la cassa scaricata al Laboratorio Alchemico. Non la apre e non scopre ancora il trasferimento sotterraneo.\\
Giorno 18 & Elias si presenta all'Archivio come verificatore delle pesature della Fonderia e chiede un confronto delle commesse. Ilyne gli mostra il registro riepilogativo e l'allegato tecnico 44-12, prima che fosse emessa la disposizione H07. Trova il marchio amministrativo e il collegamento con l'Ospedale. Nell'allegato 44-12 vede il disegno della vasca e copia la runa da incidere; Silas gli indica Miriam per interpretarla. Disegna H02; ``sotto?'' è ancora un'ipotesi.\\
Giorno 19 & Oren, suo vecchio contatto medico, lo visita e scrive H08. Conferma le forniture riservate, i brevi ritorni di lucidità seguiti da ricadute e l'aumento delle sedute malgrado le proteste. Elias teme un uso coercitivo del proprio regolatore.\\
Giorno 20 & Miriam riceve Elias su presentazione di Silas. Esamina soltanto la copia della runa. Nei testi trova associazioni a percezione, crescita e possibile contatto. Scrive H11 senza conoscere vasca, reperto o trattamento.\\
Giorno 21 & Elias ha una crisi in casa e fugge per non coinvolgere Eleanor. All'interno della chiesa arrivano voci che qualcuno sta ficcando il naso e padre Aldren viene inviato a seguire le tracce.\\
Giorni 22--24 & Elias dorme in un vano dismesso vicino alle fogne, con acqua e provviste portate da casa. Osserva il movimento delle forniture al Laboratorio.\\
Giorno 25 & Eleanor invia le lettere ai PG. Nell'Ospedale sei pazienti vengono portati nell'ala inferiore; due spariscono dai registri pubblici. Oren scrive H09.\\
Giorno 26 & Elias sfrutta un trasferimento per raggiungere il deposito e il settore merci già aperto. Ferisce al braccio il sorvegliante per liberarsi e lo risparmia; l'addetto si allontana per chiamare aiuto. Rimasto momentaneamente senza osservatori, scende con l'ascensore già abilitato, senza prenderne la chiave. Né la guardia né l'addetto assistono alla discesa. L'allarme interno parte subito e raggiunge Celestine: un intruso armato è stato visto nel deposito, poi se ne sono perse le tracce. Durante i soccorsi e i primi controlli Elias percorre la galleria, attraversa la sala del Progetto riconoscendo il regolatore e raggiunge il vano di manutenzione, dove si nasconde nell'intercapedine. Ha consultato lo schema H14 presso l'ascensore e conserva con sé le osservazioni in H12.\\
Giorno 27 & Celestine, già informata dall'allarme interno del giorno precedente, fa segnalare ai Cacciatori un intruso malato e armato di cui si sono perse le tracce nel deposito. Roland emette H10 e raccoglie la testimonianza del sorvegliante. Elias osserva le attività dal riparo nel vano di manutenzione, la curiosità per il reperto contribuisce a trattenerlo.\\
Giorno 28 & Arrivano i PG. Elias si trova nel complesso del Progetto, ancora lucido. Celestine prosegue nel suo lavoro...\\
\end{longtable}

\subsection{Elias: sabotaggio, curiosità e riparo}
Dopo la scoperta del giorno 26, Elias rimane nascosto per capire quale funzione abbia il reperto. Prepara un intervento sul sistema fra due sedute: vuole rallentarlo interrompendo la portata aumentata. Il sabotaggio potrebbe donargli il tempo di capire cosa stia accadendo davvero.

\textbf{Riparo:} Elias si nasconde in un'intercapedine bassa dietro alcune condotte. È lunga pochi metri e permette di rannicchiarsi. Elias la individua dopo l'ingresso, apre il pannello con gli attrezzi che ha portato e lo riaccosta dall'interno, un occhio attento potrebbe notarlo. Con l'avanzare del tempo tende a perdere il filo dei pensieri e riconosce la propria irritazione sempre più eccessiva.

\subsection{Che cosa fanno gli altri e perché non hanno trovato Elias}
Celestine conosce l'intrusione dall'allarme del giorno 26 e fa controllare i passaggi. Nessuno ha visto Elias prendere l'ascensore. Le ricerche riguardano il deposito e i locali e accessi praticabili del complesso. Un controllo approfondito nel vano può scoprire il pannello riaccostato. Elias si espone quando esce per osservare o agire. Al giorno 28 sono trascorsi circa due giorni dall'ingresso.

Aldren conosce alcune cose ma crede a un programma medico. Vuole recuperare le carte prima che diventino uno scandalo. Celestine conosce anche gli effetti avversi e usa la segnalazione ai Cacciatori per trovare Elias. Roland conosce i sintomi riferiti, il luogo dell'intrusione e l'aggressione: il sorvegliante ha ammesso che Elias avrebbe potuto ucciderlo e non lo ha fatto.

\section{4. Orologio della crisi}

Questo orologio misura la pressione sull'indagine. Inizia allo stadio 1 e avanza quando i giocatori perdono tempo in piste false o falliscono delle prove. Mostra i cambiamenti e permetti ai giocatori perché possano interagirci.

\begin{tabularx}{\linewidth}{>{\bfseries}p{12mm} p{35mm} X}
1 & Indagini discrete & Aldren offre aiuto e chiede delle carte. I Cacciatori verificano la segnalazione e fanno domande alla Fonderia.\\
2 & Prove confiscate & L'Archivio limita gli accessi; vengono richiesti registri e convocati testimoni.\\
3 & Vie di fuga presidiate & I Cacciatori sorvegliano le strade, un gruppo di bestie si è addentrato in città. Le persone sono più timorose, le strade diventano un rischio.\\
4 & Ricerche estese & L'ospedale viene controllato con uomini del Progetto. Entrare, uscire e fare domande diventa più difficile. I cacciatori si fanno sempre più spietati e diretti.\\
5 & Isolamento & Documenti e ingressi noti sono sotto sorveglianza. Elias comincia una trasformazione sempre meno reversibile, si avvicina il contatto con il grande essere e le persone nella città cominciano a perdere il senno anche se non sono malate, farneticano e nei casi peggiori arrivano a cavarsi gli occhi.\\
\end{tabularx}

\begin{clockbox}
Se Celestine teme che il suo piano sia compromesso può preparare lo spostamento di documentie del progetto stesso. È un'operazione visibile che i PG possono seguire o interrompere, è una mossa disperata da parte di Celestine che probabilmente porta alla fine dell'esperimento.
\end{clockbox}

\section{5. Le fazioni}

\subsection{La Chiesa della Cura}
\textbf{Volto pubblico:} ordine, guarigione, stabilità.\\
\textbf{Obiettivo:} trovare Elias prima che renda pubblico ciò che ha visto e proteggere l'istituzione, anche se alcuni vertici ignorano il Progetto.\\
\textbf{Può offrire:} lasciapassare, cure, accesso all'Ospedale, autorità, protezione per Eleanor.\\
\textbf{Prezzo:} consegnare appunti, testimonianze o Elias; accettare la supervisione di Padre Aldren.\\
\textbf{Aspetti:} \aspect{La cura giustifica il silenzio}; \aspect{Ogni porta ha un confessore}.\\
\textbf{Linea rossa:} esposizione pubblica del Progetto o accesso non autorizzato ai sotterranei.

\subsection{I Cacciatori}
\textbf{Volto pubblico:} uomini e donne armati che tengono lontane le bestie.\\
\textbf{Obiettivo:} intercettare Elias e abbatterlo so non può essere contenuto. Roland teme che la competenza tecnica renda una futura trasformazione più pericolosa.\\
\textbf{Può offrire:} scorta, uso della forza coercitiva, protezione dalle bestie e testimonianza sull'intrusione al deposito.\\
\textbf{Prezzo:} garanzie concordate per avvicinare Elias, proteggere i presenti e contenerlo se perde il controllo. Eventualmente firmare il contratto per diventare cacciatori.\\
\textbf{Aspetti:} \aspect{Ogni recupero richiede precauzioni}; \aspect{La Caccia ricorda ogni esitazione}.\\
\textbf{Linea rossa:} un'aggressione incontrollata che non può essere contenuta. Scoprire che i PG nascondono Elias provoca una richiesta di verifica e un confronto, non un abbattimento automatico.

\subsection{L'Accademia}
\textbf{Volto pubblico:} studiosi isolati, medici teorici e antiquari.\\
\textbf{Obiettivo:} capire che cosa Elias abbia scoperto e sottrarre alla Chiesa almeno una parte del materiale.\\
\textbf{Può offrire:} interpretazione delle rune, testi e analisi del materiale già disponibile.\\
\textbf{Prezzo:} richieste di campioni recuperati, se coinvolti direttamente anche l'accesso all'organismo e al corpo di Elias.\\
\textbf{Aspetti:} \aspect{Ogni limite è un'ipotesi non ancora smentita}; \aspect{La verità non promette misericordia}.\\
\textbf{Linea rossa:} distruzione completa del Progetto senza averne salvato alcuna parte.

\section{6. Personaggi: conoscenze, motivazioni e reazioni}

\begin{gmbox}
Le schede descrivono la situazione al giorno 28, prima dell'intervento dei PG.
\end{gmbox}

\subsection{Eleanor Ward}
\textbf{Ruolo e presenza:} moglie di Elias, committente; si trova a Casa Ward.\par
\textbf{Come interpretarla:} ascolta senza interrompere, poi chiede che cosa accadrà a Elias. Ha ricostruito più volte gli ultimi giorni, tiene ordinata la casa. Sembra spaventata ma decisa, non esiterà a difendere se stessa o l'onore del marito.

\textbf{Vuole, perché:} ritrovare vivo il marito nella speranza che sia lucido. Dopo risposto evasivamente alle della chiesa Chiesa, chiama persone di cui Elias si fidava. Prima della scena concorda con ogni giocatore un legame semplicecon lei o con Elias.

\textbf{Sa e fonte:}
\begin{itemize}
\item Ha assistito alle cure per la mano quattro mesi fa e ai sintomi dal giorno 12. Non conosce una diagnosi certa.
\item Elias le mostrava pesi e numeri delle spedizioni: sa che l'indagine comincia alla Fonderia. H01 e H02 sono rimasti in casa.
\item Il giorno 20 le disse di aver consultato Miriam su presentazione di Silas. Eleanor sa della visita, non delle conclusioni teoriche della studiosa.
\item Il giorno 21 vide incaricati ecclesiastici cercarlo; dopo la crisi lo vide partire con attrezzi, cibo e pistola. Non lo ha visto né ricevuto messaggi da allora.
\end{itemize}
\textbf{Crede / non sa:} pensa a forniture illecite. Non sa dove Elias si nasconda, né che ci siano esperimenti sui malati. 

\textbf{Nasconde e teme:} Elias le strinse il braccio tanto da lasciarle un livido. Lei minimizza perché teme che quel fatto basti a farlo abbattere. 

\textbf{Se ignorata:} cerca notizie alla Fonderia. Se la Chiesa offre protezione mentre i PG spariscono, può accettare per ottenere informazioni prima di mettersi alla ricerca da sola.

\subsection{Elias Ward}
\textbf{Ruolo e presenza:} Ingegnere scomparso, deciso, leale e giusto. Il suo ingegno e la sua curiosità però sono anche il suo punto debole
\par
\textbf{Come interpretarlo:} In base a quanto tempo è passato sarà più o meno lucido, ogni tanto nel mentre della conversazione zittisce i suoi interlocutori per sentire i rumori delle macchine.

\textbf{Vuole, perché:} impedire altri abusi legati alla propria invenzione e capire a cosa possa portare il progetto, questo conflitto lo trattiene.

\textbf{Sa e fonte:}
\begin{itemize}
\item Riconosce il regolatore nella macchina il giorno 26 perché l'ha progettato. Alla Fonderia aveva soltanto un sospetto. Sa della vendita perché Silas gliel'ha confessata il giorno 14.
\item Oren gli ha confermato forniture riservate e sofferenze dei pazienti. Miriam gli ha spiegato le associazioni della runa nei testi.
\item Nella sala del Progetto ha visto il regolatore incorporato nella macchina e il reperto nella vasca e i pazienti sottoposti alle sedute.
\item Elias conserva  il suo taccuino con H12 con le scoperte e i dubbi personali.
\item I sintomi precedevano l'ingresso e stanno peggiorando. Non conosce il proprio tempo residuo né una cura.
\end{itemize}

\textbf{Nasconde e teme:} quando incontra qualcuno minimizza stanchezza e spasmi, temendo di essere consegnato come cavia o considerato già una bestia. Ammette l'aggressione al sorvegliante se interrogato.

\textbf{Interazione:} racconta le scoperte verificabili e distingue ciò che ha visto dalle ipotesi di Miriam. Ammette quanto lo trattenga la curiosità.

\textbf{Sotto pressione:} se le cose si mettono male potrebbe completare la trasformazione, solo la presenza di Eleanor potrebbe calmarlo o portarlo a mettere fine alla propria vita rendendosi conto della sua trasformazione.

\textbf{Aspetti:} \aspect{Non posso lasciare un errore irrisolto}; \aspect{Le mie invenzioni hanno già versato troppo sangue}; \aspect{La bestia impara in fretta}.\par

\subsection{Padre Aldren Vale}
\textbf{Ruolo e presenza:} referente amministrativo della Chiesa. Offre l'accesso all' Archivio e all' Ospedale ma vuole controllare l'indagine.\par
\textbf{Come interpretarlo:} prende nota delle richieste, offre una sedia e trasforma ogni urgenza in una procedura. Cerca di essere amico di tutti. Chiede di vedere i documenti prima di commentarli. Legale neutrale purissimo.

\textbf{Vuole, perché:} recuperare le carte e ricondurre Elias sotto supervisione. Ritiene che la fiducia nella Chiesa sia fonadamentale per la città e teme che uno scandalo danneggi anche le cure che funzionano. Protegge inoltre la propria posizione. Può essere convinto ad accusare Celestine se avesse le prove in mano.

\textbf{Sa e fonte:}
\begin{itemize}
\item Sa che le forniture scendono nell'ala inferiore e che Celestine dirige il programma; non è entrato nella sala del Progetto. Può chiedere un incontro ma non è detto che gli venga assicurato.
\end{itemize}
\textbf{Crede / non sa:} Crede che Celestine diriga una ricerca clinica rischiosa ma giustificabile. Non conosce i trattamenti imposti senza consenso, le sperimentazioni sui sei pazienti, gli effetti avversi occultati o la natura occulta del Progetto.

\textbf{Nasconde e teme:} Tiene le informazioni per se fino a che non capisce che Celestine sia un pericolo per la chiesa. Teme che i PG mettano in cattiva luce la chiesa.

\textbf{Interazione:} rivela che esiste una pratica di ricerca e offre consultazione limitata o accesso all'archivio. Se gli mostrano H09 e gli permettono di confrontare i registri, può iniziare a dubitare.

\subsection{Capitana Roland Kreiss}
\textbf{Ruolo e presenza:} comandante dei Cacciatori;
Può fornire forza, contenimento e scorta.\par
\textbf{Come interpretarla:} cacciatore di esperienza, crede solo a quello che vede, è ai limiti del paranoico. Come tutti i cacciatori ha firmato col sangue il patto e vive l'incubo del cacciatore.

\textbf{Vuole, perché:} evitare vittime. In passato esitò ad abbattere una persona già trasformata che sembrava riconoscerla; la successiva aggressione costò cinque vite. Da allora richiede prove di autocontrollo agli stessi cacciatori.

\textbf{Sa e fonte:}
\begin{itemize}
\item Il giorno 27 ricevette dall'Ospedale una segnalazione di intruso malato e armato. Il contatto fu organizzato da Celestine, ma Roland conosce solo la provenienza ospedaliera.
\item Il sorvegliante ferito le raccontò che Elias lo aveva colpito al braccio per liberarsi e poi risparmiato.
\item Conosce il luogo dell'intrusione, nel deposito interno al Laboratorio Alchemico, dalla segnalazione e dal racconto del sorvegliante.
\end{itemize}
\textbf{Crede / non sa:} la competenza tecnica renderebbe Elias pericoloso se si dovesse trasformare. Non conosce l'esperimento e non sa che la denuncia ometta deliberatamente il contesto. Non conosce l'accesso tramite le fogne.

\textbf{Interazione:} mostra H10 e riferisce luogo dell'intrusione e testimonianza del ferito. Chiede prove della lucidità di Elias (H08 prova lucidità al giorno 19, non quella attuale). Se messa al corrente della pericolosità del piano potrebbe scatenare i suoi cacciatori contro Celestine.

\textbf{Se ingannata:} rafforza le pattuglie, chiede a Nox di seguire i PG. Può rivedere la valutazione davanti a prove nuove. \par

\subsection{Dottoressa Miriam Voss}
\textbf{Ruolo e presenza:} studiosa dei testi pthumeriani, si trova nella Biblioteca Proibita, indipendente ma legata agli studiosi dell'Accademia. Silas la conosce per consulenze su testi e strumenti e le ha presentato Elias.\par
\textbf{Come interpretarla:} studiosa, molto tesa e con mille pensieri in testa, a volte sembra quasi parlare da sola. 

\textbf{Vuole, perché:} comprendere l'impiego moderno ddella runa che Elias gli ha mostrato. Teme che testimonianze e documenti vengano distrutti prima di poterli studiare.

\textbf{Sa e fonte:} il giorno 20 Elias le mostrò soltanto la copia isolata della runa trovata all'Archivio. Nei suoi testi la runa alcuni passi sembrano collegarla alla possibilità di contatto con esseri superiori. H11 contiene questa lettura. Può mostrarne una fonte in H11b, preghiera tradotta e commentata da Willem.

\textbf{Non sa:} ignora vasca, reperto, pazienti, regolatore e ruolo di Celestine. Non identifica la ricerca dal marchio amministrativo.

\textbf{Interazione:} spiega la runa e racconta il colloquio senza esigere campioni come pedaggio. Se vede H06b riconosce che l'incisione era destinata a una vasca, ma non sa che cosa contenga. Se riceve anche referti e H16 può confrontare la teoria di Celestine con i risultati.

\textbf{Aiuto e limite:} traduce ed aiuta a confrontare testi e prove. Se lasciata libera di studiare può arrivare a capire che c'è il rischio di contatto con un grande essere e che questo possa essere devastante ma il modo di esprimerlo potrebbe farla passare per pazza.

\subsection{Madre Celestine Vey}
\textbf{Ruolo e presenza:} responsabile della ricerca nell'ala inferiore dell'Ospedale. Dirige le prove e le forniture riservate.\par
\textbf{Come interpretarla:} agisce con autorità fascista nascosta da un fare paternalista e materno. Vuole astrarsi dal suo corpo mortale entrando in connessione con i grandi esseri che non conosce davvero fino in fondo, consuma grandi quantitativi di sangue antico cercando di elevarsi. 

\textbf{Vuole, perché:} rendere duraturi i miglioramenti che ha davvero osservato. Pensa che fermarsi significhi rinunciare a una cura ed una stabilizzazione nell'uso del sangue. Aveva iniziato convinta che i suoi sforzi fossero per un bene comune ma non è più in grado di distinguere il bene o il male. Tuttavia rimane estremamente lucida e capace di qualsiasi interazione sociale, per ora.

\textbf{Sa e fonte:} ha seguito la ricerca fin dalle prime prove volontarie. I ritorni di lucidità svaniscono; i pazienti ricadono nello stato precedente. Sa che ora le sedute vengono imposte e che i rapporti pubblici minimizzano ricadute ed effetti avversi. Ha ordinato l'incisione della runa e l'aggiunta del regolatore per aumentare la portata. L'allarme interno del giorno 26 l'ha informata dell'intrusione di Elias nel deposito. Nessun testimone lo ha visto scendere: sospetta che possa aver raggiunto l'ala inferiore, ma non ne ha conferma e non conosce il nascondiglio.

\textbf{Crede / non sa:} attribuisce il beneficio al Grande Essere a cui ritiene appartenesse il reperto. Considera ogni ricaduta un segno di contatto insufficiente. Non ha verificato provenienza, efficacia della runa o comunicazione con un'entità.

\textbf{Nasconde e teme:} occulta proteste, rifiuti e fallimento del recupero duraturo. Teme che una sospensione renda inutili i sacrifici già compiuti; usa questo timore per giustificarne altri.

\textbf{Sotto minaccia:} cerca di fermare gli intrusi e preservare reperto e documenti. Davanti a un assalto imminente può preparare un trasferimento visibile, se ha una via praticabile.

\subsection{Dottor Oren Sile}
\textbf{Ruolo e presenza:} medico dell'Ospedale Vecchio, testimone degli effetti e possibile soccorritore. Non fa parte del Progetto.\par
\textbf{Come interpretarlo:} abbassa la voce quando parla di cose pericolose, ma torna preciso quando descrive un sintomo. Tocca insistentemente i suoi occhiali sottili, non riesce a capire chi siano i buoni e chi i cattivi. La competenza clinica rimane anche quando il coraggio gli manca. 

\textbf{Vuole, perché:} evitare altre vittime e sottrarsi alle conseguenze della denuncia. Vorrebbe fuggire, ma lasciare i malati confermerebbe l'idea di essere diventato complice.

\textbf{Sa e fonte:} conosce Elias da precedenti collaborazioni mediche. Il giorno 19 lo visita e gli conferma le forniture riservate. Ha seguito gli effetti clinici dei primi trattamenti volontari, osservando miglioramenti seguiti da ricadute; vede la vasca soltanto il giorno 25. Quel giorno accompagna i sei lungo la scala interna fino all'ala inferiore, malgrado il loro rifiuto; nella sala del Progetto vede i collegamenti alla vasca e scrive H09. Conosce il percorso ma non può autorizzare estranei a scendere. Non conosce la finalità ultraterrena.

\textbf{Nasconde e teme:} tace inizialmente sull' avere preparato un trasferimento senza consenso. Teme arresto, perdita della professione e disprezzo. Conserva H09 senza spedirlo perché non sa a chi affidarlo senza esporre un altro testimone.

\textbf{Interazione:} fornisce visita e H08 a chi sta cercando Elias; una conversazione privata evita di esporlo al personale ecclesiastico. Può essere una fonte di padre Aldren o altri membri della chiesa a descrizione del GM.

\textbf{Se è stato interrogato dalla chiesa:} prepara la fuga, abbandona H09 nel suo ufficio. Minacciarlo può ottenere una confessione.\par
\textbf{Voce:} ``Pensavo che un superiore avrebbe fermato un errore. Gli ho consegnato l'uomo che se n'era accorto.''

\subsection{Silas ``Tre Dita''}
\textbf{Ruolo e presenza:} ingegnere e fabbro all'Officina delle Lame; collega il passato di Elias alle forniture attuali.\par
\textbf{Come interpretarlo:} continua a lavorare mentre parla di affari. Perse le sue dita lavorando distrattamente, preda dei vizi. Usa l'ironia per coprire le proprie debolezze

\textbf{Vuole, perché:} aiutare Elias e non sentirsi un traditore. Vendette la copia degli schemi del regolatore per pagare i propri debiti, volendo credere una giustificazione ospedaliera che gli permetteva di non interrogarsi troppo sull'utilizzo del regolatore.

\textbf{Sa e fonte:} lavorava con Elias, conosce l'impiego originario del regolatore e il motivo dell'incendio. Vendette la copia sette mesi fa e riconosce il marchio amministrativo del compratore, presentatosi come incaricato della Chiesa; non conosceva il nome di Celestine. Una commessa successiva al Laboratorio gli diede una chiave dell'ingresso di servizio che dà sul deposito interno, rimasta in suo possesso. Apre solo quell'ingresso. H03 è un frammento originale di Elias conservato da Silas. Il giorno 14 confessò la vendita a Elias; in seguito gli procurò l'incontro con Miriam, che conosce per consulenze su vecchi strumenti.

\textbf{Nasconde e teme:} Teme la rabbia di Elias e un'accusa di traffico di materiale ecclesiastico.

\textbf{Interazione:} mostra H03 e spiega il vecchio progetto senza pretendere fiducia assoluta. Il nome «regolatore» su H05 è compatibile con l'invenzione, ma non basta a identificarla.

\textbf{Aiuto e limite:} consegna la chiave se i PG la chiedono. Può fare da testimone riconoscendo un componente noto o testimoniare sulla vendita. Se gli viene dato abbastanza tempo può perfino fornire strumenti come armi o gadget particolari.

\textbf{Se ignorato:} Se la Chiesa reagisce alle indagini può arrestarlo per paura di confessioni verso i PG. Minacciato dai PG, chiude bottega e cerca protezione.
\par
\textbf{Voce:} ``Mi dissero che serviva per un ospedale. Era vero. È tutto il resto che non ho voluto chiedere.''

\subsection{Marta Ghal}
\textbf{Ruolo e interpretazione:} caporeparto della Fonderia; brusca, controlla che le conversazioni siano confidenziali. Vuole proteggere gli operai e il loro salario, incluso Elias.

\textbf{Sa e fonte:} conserva il foglio originale H04. Il giorno 16 registrò Elias come presente mentre lui seguiva il carro di T. Harrow con la cassa 17-B. Sapeva che voleva verificarne la destinazione, ma non lo accompagnò.

\textbf{Crede / nasconde / teme:} sospetta frodi nelle commesse. Nasconde la falsa registrazione della presenza di Elias: se scoperta, potrebbe costarle il posto.

\textbf{Reazione:} se scatta la confisca, conserva o consegna una copia quando ha tempo materiale per farlo; altrimenti ricorda numeri e testimoni. Una minaccia ai lavoratori la rende ostile anche se i PG cercano Elias.\par
\textbf{Voce:} ``Il peso l'abbiamo controllato. Chiedete a chi ha cambiato il registro, non a chi reggeva la bilancia.''

\subsection{Sorella Ilyne}
\textbf{Ruolo e interpretazione:} archivista ecclesiastica; controlla firme e numeri prima di dare un'opinione. Crede che registrazioni corrette permettano alla Chiesa di rispondere dei propri atti. Donna anziana, alta e magra, voce incerta, chiede sempre conferma di aver capito bene. 

\textbf{Sa e fonte:} gestisce il registro H06, la scheda 44-14 (H06a) e il disegno 44-12 (H06b), tenendoli distinti dalle disposizioni interne. Il giorno 18 mostrò a Elias, identificatosi come verificatore delle pesature, il riepilogo e l'allegato 44-12. Custodisce H07 fra le disposizioni interne nel deposito perché deve applicare l'ordine relativo a Elias. Riconosce il marchio delle commesse riservate ma non sa chi diriga la ricerca.

\textbf{Crede / nasconde / teme:} sospetta che qualcuno usi l'autorità ecclesiastica senza risponderne. Non offre subito il proprio sospetto perché teme di accusare ingiustamente colleghi e perdere l'accesso all'Archivio.

\textbf{Interazione e aiuto:} il numero di commessa individua la pratica. Porta il registro riepilogativo e gli allegati ordinari, senza ammettere i visitatori nel deposito. La scheda 44-14 e il disegno 44-12 sono consultabili per una verifica motivata. Quando i PG dichiarano di cercare Elias, segnala l'esistenza di una disposizione successiva alla sua visita e rimanda ad Aldren per l'autorizzazione. Può mostrare H07 in privato se minacciata o convinta.

\subsection{Maud Kerren}
\textbf{Ruolo e interpretazione:} infermiera dell'Ospedale; ricorda abitudini e bisogni dei degenti, parla dei pazienti per nome. Vuole evitarne un altro trasferimento senza spiegazioni.

\textbf{Sa e fonte:} assisteva i sei prima del trasferimento del giorno 25 e può riconoscerli. Ha visto i brevi ritorni di lucidità e le ricadute; ha raccolto le richieste dei pazienti di non continuare le sedute.

\textbf{Crede / nasconde / teme:} sospetta esperimenti dannosi, ma non sa che cosa misurino. Teme che perdere il posto significhi perdere ogni possibilità di seguire i malati.

\textbf{Interazione e aiuto:} racconta il trasferimento e mostra la lista a chi cerca i malati. Conosce H09 solo se Oren glielo affida o qualcuno lo trova. Può guidare nel reparto e assistere i pazienti.

\textbf{Reazione:}Può chiedere l'aiuto dei PG ed accompagnarli se riceve protezione per cercare di salvare i malati.

\subsection{Nox}
\textbf{Ruolo e interpretazione:} giovane cacciatore agli ordini di Roland; controlla spesso la reazione della capitana prima di parlare. Vuole essere considerato affidabile e teme che prudenza significhi codardia.

\textbf{Sa e fonte:} conosce H10 e gli ordini ricevuti; sa solo ciò che ha visto durante le pattuglie.

\textbf{Crede / nasconde / teme:} interpreta l'ansia della capitana come invito ad anticipare il pericolo. Nasconde i propri dubbi ostentando sicurezza.

\textbf{Interazione e aiuto:} può guidare una pattuglia e aiutare i PG se si dovessero trovare in pericolo. I PG possono convincerlo della pericolosità del Progetto e chiedergli aiuto militare.

\subsection{Abel Norn}
\textbf{Ruolo e interpretazione:} bibliotecario cieco; conosce l'ordine dei volumi e le abitudini degli studiosi.

\textbf{Sa e fonte:} riconobbe Miriam accompagnata da un visitatore il giorno 20 e ne udì la presentazione come Elias Ward. Può ritrovare, con i contrassegni degli scaffali o l'aiuto di un assistente, il testo consultato. Non vide la copia della runa e non lesse gli appunti privati di Elias.

\textbf{Aiuto e reazione:} se Miriam manca, può far consultare il volume con la traduzione e i commenti di Willem riprodotti in H11b.

\subsection{Personale operativo: risposte per il master}
Queste figure completano le azioni già presenti nella trama. 

\textbf{T. Harrow, vetturale di H05.} Ha trasportato il collo 17-B al deposito il giorno 16 e conosce la ricodifica perché ha firmato la ricevuta. Conferma percorso, peso e firma; non conosce il contenuto oltre la descrizione e non ha consegnato la cassa nella sala del Progetto. Un colloquio privato o una copia di H05 basta a rendere concreta la domanda.

\textbf{Contabile ecclesiastico alla Fonderia.} Ha applicato la descrizione di copertura e autorizzato 17-B. Conosce il carattere riservato delle forniture, non il contatto o le sperimentazioni sui sei pazienti. Difende inizialmente la ricodifica come procedura; davanti ai due documenti deve ammettere che il peso standard non fu misurato. Teme un'accusa di frode e segnala ai superiori chi ha visto le carte. Può confiscare registri tramite l'autorità ecclesiastica, ma non eliminare copie di cui ignora l'esistenza.

\textbf{Sorvegliante ferito.} È l'uomo colpito da Elias il giorno 26. Presidiava il trasferimento e tentò di trattenerlo. Ha riferito a Celestine e poi a Roland che Elias lo colpì al braccio e fuggì senza finirlo. Vuole cure e non essere accusato di aver lasciato passare l'intruso. Una domanda sul gesto concreto può confermare che fu risparmiato; non può certificare che Elias rimarrà lucido. Può essere sentito attraverso Roland o rintracciato al posto di cura dell'Ospedale. Non è contemporaneamente la guardia in servizio al deposito.

\textbf{Due tecnici del Laboratorio.} L'addetto all'ascensore conosce inventario H13, carichi e galleria, e custodisce la chiave che abilita il mezzo. Il giorno 26 si allontanò per chiamare aiuto dopo che Elias ebbe ferito il sorvegliante. Aveva lasciato l'ascensore abilitato per il carico, ma non vide Elias usarlo. Il tecnico della macchina conosce la funzione del regolatore e si trova al deposito quando il lavoro su forniture o manutenzione lo richiede; non sa come ottenere un contatto. Entrambi temono punizioni per accessi non autorizzati.

\textbf{Sorvegliante in servizio al deposito.} È una guardia diversa dal ferito. Controlla lasciapassare e passaggi riservati. Sa dell'intruso e dell'obbligo di riferire, ma non possiede tutte le conoscenze di Celestine.

\section{7. Mappa della rete}

\begin{center}
\resizebox{\linewidth}{!}{%
\begin{tikzpicture}[
 node distance=13mm and 17mm,
 every node/.style={font=\small,align=center},
 main/.style={draw=blood,thick,fill=ash!55,minimum width=31mm,minimum height=10mm},
 nodebox/.style={draw=iron,fill=mist,minimum width=30mm,minimum height=9mm},
 final/.style={draw=blood,very thick,fill=blood!10,minimum width=34mm,minimum height=11mm},
 edge/.style={-{Latex[length=2mm]},thin,iron}
]
\node[main] (home) {Casa Ward};
\node[nodebox,below left=of home] (foundry) {Fonderia};
\node[nodebox,below=of home] (studio) {Studio distrutto - Officina delle lame};
\node[nodebox,below right=of home] (hunters) {Rifugio dei Cacciatori};
\node[nodebox,below left=of foundry] (archive) {Archivio Ecclesiastico};
\node[nodebox,below=of studio] (hospital) {Ospedale Vecchio};
\node[nodebox,below right=of hunters] (library) {Biblioteca Proibita};
\node[nodebox,below=19mm of hospital] (alchemy) {Laboratorio Alchemico};
\node[final,below=17mm of alchemy] (project) {Il Progetto};
\draw[edge] (home)--(foundry);
\draw[edge] (home)--(studio);
\draw[edge] (home)--(hunters);
\draw[edge] (foundry)--(archive);
\draw[edge] (foundry)--(alchemy);
\draw[edge] (studio)--(archive);
\draw[edge] (studio)--(library);
\draw[edge] (hunters)--(hospital);
\draw[edge] (hunters)--(alchemy);
\draw[edge] (archive)--(hospital);
\draw[edge,bend left=20] (archive) to (library);
\draw[edge] (hospital)--(alchemy);
\draw[edge] (library)--(alchemy);
\draw[edge] (alchemy)--(project);
\draw[edge,bend right=20] (hospital) to (project);
\end{tikzpicture}%
}
\end{center}

Le frecce indicano collegamenti disponibili. I PG possono raggiungere il nodo finale quando hanno individuato e superato almeno un accesso fisico all'ala inferiore, anche direttamente dall'Ospedale.

\medskip

\section{8. Matrice delle rivelazioni}

\small
\begin{longtable}{p{37mm} p{34mm} p{34mm} p{43mm}}
\toprule
\textbf{Rivelazione} & \textbf{Prima fonte} & \textbf{Alternativa} & \textbf{Conferma / limite}\\
\midrule
Elias indagava & Casa: H01, H02 & Marta alla Fonderia & Silas ricorda la visita e la confessione\\
Origine del regolatore & Officina: H03, Silas & Elias al finale & Fonderia prova la presenza, non da sola la paternità\\
Spedizione occultata & Fonderia: H04 e H05 & Marta e vetturale & Archivio: H06a conferma la spedizione; H07 aggiunge le disposizioni riservate\\
Ospedale coinvolto & Archivio: H06 & Ospedale: Oren, H09 & Laboratorio: H13 registra consegne dirette\\
Deposito di transito & Laboratorio: H13 & Tecnici e ascensore & Fonderia indica la destinazione, non il tratto successivo\\
Accesso al complesso & Laboratorio: H14 & Scala sorvegliata dell'Ospedale & Addetti e rimando in H13; H12 non necessario\\
Restrizioni e ricerca & Archivio: H07, Ilyne & Oren e Aldren & Finale: H16 mostra ciò che Aldren ignorava\\
Lettura della runa e obiettivo & Biblioteca: H11 e H11b, lettura e fonte & Finale: H16, dichiarazione di Celestine & H15 mostra il circuito del trattamento, non prova un contatto\\
Autocontrollo di Elias & Rifugio: sorvegliante risparmiato & Finale: incontro con Elias & H08 prova lucidità al giorno 19; la competenza tecnica non basta\\
\bottomrule
\end{longtable}
\normalsize

\nodehead{9. Nodo 0 -- Casa Ward}{Innesco, legame emotivo e primo ventaglio di piste.}{Elias è fuggito, è stato rapito o stava facendo qualcosa di preciso?}

\subsection{A colpo d'occhio}
\begin{overviewbox}
\begin{tabularx}{\linewidth}{>{\bfseries}p{32mm} X}
Presenti & Eleanor; Aldren soltanto per un arrivo motivato.\\
Accesso & Eleanor riceve i PG in casa e mette a disposizione le carte rimaste.\\
Handout & H01 e H02 sul tavolino dello studio domestico.\\
Da ricordare & Elias è partito volontariamente il giorno 21. Eleanor non conosce il suo nascondiglio.\\
\end{tabularx}
\end{overviewbox}

\subsubsection{Situazione iniziale}
Elias è uscito volontariamente il giorno 21, sette giorni prima dell'arrivo dei PG. Aveva visto incaricati ecclesiastici cercarlo e, durante una crisi, aveva stretto il braccio di Eleanor. Spaventato da sé stesso, ha preso attrezzi, pistola, provviste e il taccuino recente. Gli appunti precedenti sono rimasti in casa. Nessuno lo ha rapito qui e non c'è stata una perquisizione segreta.

Eleanor ha continuato a vivere nella casa e a cercare notizie. Dopo i tentativi con Guardia e Chiesa ha scritto ai PG il giorno 25. Oggi ha raccolto sul tavolino le carte rimaste, perché possano esaminarle. La scena deve dare un motivo umano alla ricerca e più piste praticabili.

\subsection{Arrivo e ambienti}
\subsubsection{Ingresso}
\begin{readaloud}
Eleanor apre e vi fa entrare senza trattenervi sulla soglia. Nella stanza c'è odore di carbone e di minestra. Il tavolo è pulito; sopra una sedia è piegato un cappotto da lavoro pesante.

Attraverso una porta aperta vedete un tavolino coperto di fogli, un astuccio di strumenti e una cartina ripiegata. Eleanor sposta una tazza per farvi posto. Quando tira indietro la sedia si aggiusta la manica sul polso.

``Queste sono le carte che ha lasciato. Ho chiesto alla Guardia, poi alla Chiesa. Nessuno mi ha detto di averlo trovato. Da dove volete cominciare?''
\end{readaloud}

\begin{gmbox}
la casa è stata abitata e riordinata durante la settimana. Minestra e tazza appartengono alla giornata corrente. Il cappotto è un ricambio, non una prova di ritorno. La manica copre il livido della crisi. L'ordine domestico mostra lo sforzo di Eleanor di continuare; non nasconde un enigma.
\end{gmbox}

\textbf{Aspetti di scena:} \aspect{Una casa rimasta in attesa}; \aspect{Ogni oggetto ricorda una vita abbandonata}.

\subsubsection{Spazi ed elementi esaminabili}
La stanza principale comunica con un piccolo studio domestico. Questo studio non è il laboratorio incendiato due anni fa, che si visita con l'Officina delle Lame. La camera da letto e la dispensa contengono oggetti ordinari: Eleanor può spiegare che cosa Elias abbia portato via senza richiedere una perquisizione completa.

\begin{description}[leftmargin=5mm,style=nextline]
\item[Carte sul tavolino: H01 e H02.]
Eleanor le ha raccolte dai fogli di lavoro lasciati da Elias; H02 è ripiegato accanto a H01. Le schede nella sezione Handout riportano consegna e interpretazione.
\item[Strumenti e vecchi disegni.]
Restano compassi, strumenti di misura e schizzi di pezzi isolati. Eleanor spiega che Elias li conservava dal lavoro d'ingegnere, mentre portò via gli attrezzi da intervento. I disegni non contengono una macchina completa e non permettono di ricostruire il Progetto. Per identificarne l'origine si può consultare Silas; qui non si trova H03.
\item[Il materiale che manca.]
Eleanor ha visto Elias prendere borsa, pistola, cibo e taccuino. È una testimonianza di partenza preparata in fretta.
\item[Il polso di Eleanor.]
Se i PG lo notano, descrivi un livido in via di guarigione. Lei si copre quando la conversazione riguarda l'aggressività di Elias. Il segno sostiene il suo racconto della crisi; non stabilisce una diagnosi. Se non viene notato, la domanda sulle condizioni di Elias può comunque portare alla confessione.
\end{description}

\subsection{Interazioni e ostacoli}
\subsubsection{Parlare con Eleanor}
La scheda di Eleanor nella sezione 6 ne descrive paure e limiti. Qui è presente da sola, salvo un arrivo motivato di Aldren. Non conosce i sei pazienti, l'organismo o il nascondiglio di Elias.

\textbf{Racconta senza pressione:} l'assenza dura dal giorno 21; Elias lavorava alla Fonderia, controllava pesi e commesse e aveva ripreso a parlare dei vecchi progetti. Aveva ricevuto sangue per la mano e nelle ultime settimane dormiva male. Dice quali oggetti portò con sé e mostra le carte.

\begin{description}[leftmargin=5mm,style=nextline]
\item[``Perché pensi che non sia semplicemente scappato?'']
``L'ho visto andare via. Ma prima cercava qualcosa in quei registri, e c'era gente della Chiesa che lo cercava.'' Conferma che partì volontariamente; non sa se in seguito sia stato catturato. Non trasformare la sua speranza in una certezza investigativa.
\item[``Che cosa lo preoccupava al lavoro?'']
«Gli avevano chiesto di costruire il suo regolatore e aveva rifiutato. Poi aveva notato una cassa per lo stesso committente, con un peso diverso da quello scritto. Temeva che avessero trovato qualcun altro.» Eleanor riferisce ciò che Elias le aveva raccontato all'inizio dell'indagine. Non sa che cosa ci fosse nella cassa né conosce la confessione di Silas.
\item[``Era malato? Ti ha fatto male?'']
Ammette insonnia e irritabilità. Sulla crisi esita: ``Aveva paura. Non voglio che questo diventi una condanna.'' Se i PG spiegano come useranno l'informazione per preparare un recupero sicuro, racconta la stretta al braccio e la fuga immediata. Non pretende che promettano di ignorare un'aggressione futura.
\item[``Chi aveva consultato?'']
Ricorda Silas e il medico Oren Sile, vecchi contatti del marito. Il giorno 20 Elias le disse di essere stato da Miriam, su presentazione di Silas. Può indicare dove lavorano Silas e Oren; per raggiungere Miriam suggerisce di chiedere a Silas.
\item[«Che cosa gli disse Miriam?»]
«Non me lo spiegò. Aveva copiato un segno che non capiva. Silas lo aveva indirizzato da lei.» Eleanor può indicare il contatto, non anticipare le teorie della Biblioteca.
\item[``Possiamo portare via le carte?'']
Acconsente, chiedendo che le vengano restituite o copiate. Non condiziona le piste a un favore. Se si parla di consegnarle alla Chiesa, vuole sapere chi le custodirà e se i PG conserveranno una copia.
\end{description}

\subsubsection{Intervento eventuale di Padre Aldren}
Fallo arrivare solo se Eleanor lo ha contattato, i PG hanno chiesto un incontro, oppure sta eseguendo l'offerta di aiuto prevista allo stadio 1 dell'orologio. In quest'ultimo caso conosce l'indirizzo dalla pratica su Elias: viene a verificare se sia tornato e a chiedere le carte rimaste, non perché sappia automaticamente che i PG le stanno leggendo.

Bussa, si presenta e chiede di parlare. Offre di controllare la pratica o facilitare una visita al reparto pubblico. Domanda quali documenti siano stati trovati e propone di prenderli in custodia. Non conosce il contenuto di fogli che non ha visto. Non rivela il Progetto né ordina un arresto per la sola presenza dei PG.

Se i PG accettano, registra esattamente che cosa gli consegnano o raccontano e quale accesso promette in cambio. Se rifiutano, può lasciare un recapito e riferire il rifiuto.

\subsubsection{Azioni e complicazioni}
\begin{fatebox}
\textbf{Nessun tiro per:} ricevere H01/H02, leggere i pesi, vedere i luoghi segnati o conoscere le visite che Eleanor ricorda. Una domanda pertinente basta.

\textbf{Empatia o Carisma:} utili se i PG cercano di affrontare la crisi dopo che Eleanor si è chiusa. Un fallimento può lasciare la confessione incompleta o farle chiedere una garanzia concreta; non cancella le piste. Se l'approccio risolve già la sua paura, non chiedere un tiro per formalità.

\textbf{Indagine o Tecnica:} possono creare un vantaggio nel confronto dei documenti o distinguere disegni vecchi e appunti recenti.

\textbf{Vantaggi possibili:} \aspect{Eleanor ci ha raccontato la crisi}; \aspect{Copie delle carte al sicuro}; \aspect{Un incontro concordato con Oren}. Quest'ultimo richiede che il contatto sia effettivamente avvenuto, non soltanto un tiro riuscito.

Usa eventuali Compel soltanto su legami e Aspetti stabiliti con i giocatori. Non inventare un debito con Aldren o un vincolo familiare per imporre la collaborazione.
\end{fatebox}

\subsection{Handout}
\label{nodo-0-handout}
\phantomsection
\addcontentsline{toc}{subsubsection}{H01 -- Appunti sparsi}
\begin{handoutbox}{H01}{Appunti sparsi}
\textbf{Dove:} Sul tavolino dello studio domestico, raccolti da Eleanor. È un foglio lasciato da Elias, distinto dal taccuino che porta con sé.

\textbf{Quando consegnarlo:} Quando i PG guardano le carte o chiedono dell'indagine; nessun tiro.

\textbf{Che cosa rivela:} Il peso dichiarato è 412 libbre; quello rilevato è 391 sulla bilancia grande e 391 e un quarto sulla piccola. L'arrotondamento è compatibile con due strumenti di diversa precisione. La richiesta rifiutata e il committente comune spiegano perché Elias abbia controllato la spedizione; il peso diverso lo spinge a verificare ordine e destinazione nei registri della Fonderia.

\textbf{Limiti:} Elias non ha visto il contenuto della cassa. Il peso inferiore è un'anomalia da verificare, non la prova immediata di un furto. H05 dichiara un regolatore, ma la paternità del progetto richiede ancora il confronto con Silas e H03. ``Non parlarne con Marta'' precede il favore sul turno del giorno 16 e non dimostra che Marta sia complice.
\end{handoutbox}

\phantomsection
\addcontentsline{toc}{subsubsection}{H02 -- Cartina annotata}
\begin{handoutbox}{H02}{Cartina annotata}
\textbf{Dove:} Ripiegata sullo stesso tavolino, accanto a H01.

\textbf{Quando consegnarlo:} Insieme a H01, appena viene aperta; nessun tiro.

\textbf{Che cosa rivela:} Fonderia, Archivio e Ospedale sono cerchiati; è indicato anche il Laboratorio Alchemico, collegato da frecce. Tutti sono visitabili: i PG non devono scegliere per forza la Fonderia.

\textbf{Limiti:} Le frecce sono collegamenti studiati da Elias, non l'ordine obbligatorio delle visite. ``Sotto?'' è una domanda, non la posizione di un ingresso. ``Non segue la strada'' esprime il sospetto che il tragitto visibile non spieghi la destinazione finale; la mappa non disegna un passaggio percorribile. L'appunto sull'Archivio indica dove verificare le registrazioni, non identifica l'autore materiale delle falsificazioni.
\end{handoutbox}

\textbf{Se i PG non esaminano le carte:} Eleanor offre di mostrarle prima della partenza, senza insistere su una destinazione. Se rifiutano, può comunque raccontare l'anomalia e indicare Fonderia, Silas e Oren. Non ricorda ogni numero o annotazione: le prove rimangono in casa finché qualcuno le sposta. Se i PG scelgono una pista diversa, la scena ha comunque funzionato.

\subsection{Evoluzione della scena}
\begin{description}[leftmargin=5mm,style=nextline]
\item[Collaborano e danno notizie.]
Eleanor rimane un contatto affidabile, conserva quanto le affidano e può riconoscere la grafia in prove successive. Un piano di protezione concordato riduce la necessità di rivolgersi alla Chiesa.
\item[La minacciano o frugano contro la sua volontà.]
Chiede loro di uscire e può cercare aiuto da Marta. I documenti già letti non smettono di esistere; il costo è la fiducia e l'eventuale intervento di chi viene chiamato, non la cancellazione dell'indagine.
\item[Attirano pubblicamente l'attenzione.]
Se accusano la Chiesa in strada o coinvolgono le autorità, una segnalazione può raggiungere Aldren. Stabilisci chi ha ascoltato e che cosa riferisce. La casa non diventa sorvegliata per il semplice passare della scena.
\item[Partono o scelgono di non indagare qui.]
Non introdurre un assalto per trattenerli. Eleanor chiede quando potrà avere notizie. Se resta a lungo senza risposte cerca Marta e può accettare un successivo aiuto ecclesiastico; lascia ai PG un messaggio o un testimone del suo spostamento.
\end{description}

\begin{clockbox}
Una conversazione sensata non fa avanzare da sola l'orologio. Se passa una fase della giornata, applica la progressione generale della sorveglianza. L'offerta di protezione alla casa avviene solo per un'azione concreta di Aldren o una richiesta di Eleanor. Se la accetta, gli incaricati controllano gli ingressi e chiedono di visionare le carte: descrivi chi arriva, che cosa prende e dove lo porta.
\end{clockbox}

\textbf{Prima di chiudere la scena:} annota quali carte sono state prese o copiate, che cosa è stato detto ad Aldren, se Eleanor ha raccontato la crisi e quale pista i PG intendono seguire. Per partire è sufficiente una domanda verificabile su Elias; la verità dell'esperimento emergerà altrove.

\subsection{Collegamenti}
\begin{longtable}{p{34mm} p{60mm} p{54mm}}
\toprule
\textbf{Destinazione} & \textbf{Che cosa la suggerisce} & \textbf{Domanda da verificare}\\
\midrule
Fonderia & H01, mestiere di Elias e nome di Marta. & Il peso è confermato dai registri? Che cosa è partito?\\
Officina delle Lame & Vecchi disegni e ricordo di Silas. & Chi può riconoscere il regolatore e ricostruirne la storia?\\
Archivio & Annotazioni di H01 e H02. & Quali commesse e destinatari corrispondono alle spedizioni?\\
Ospedale Vecchio & H02, cure di Elias e nome di Oren. & Che cosa osservò il medico e perché Elias andò da lui?\\
Laboratorio Alchemico & Nome e frecce in H02. & È davvero la destinazione dei componenti?\\
Biblioteca, tramite Silas & Racconto del colloquio con Miriam. & Che cosa significava il segno copiato da Elias?\\
Rifugio dei Cacciatori & Decisione dei PG di consultare chi tratta le trasformazioni. & Esiste già una segnalazione su Elias? Eleanor non conosce H10.\\
\bottomrule
\end{longtable}

\nodehead{10. Nodo 1 -- La Fonderia}{Nodo tecnico e sociale; verifica i sospetti di Elias e ricostruisce una spedizione.}{Che cosa è partito, per conto di chi e verso quale destinazione?}

\subsection{A colpo d'occhio}
\begin{overviewbox}
\begin{tabularx}{\linewidth}{>{\bfseries}p{32mm} X}
Presenti & Marta, operai e contabile della Chiesa; Harrow e Nox secondo le condizioni descritte nelle interazioni.\\
Accesso & Marta riceve le domande in ufficio; chiede di consultare gli originali sul posto.\\
Handout & H04 nell'ufficio di Marta; H05 fra le copie carbone del magazzino.\\
Da ricordare & La cassa 17 è partita il giorno 16 come 17-B. Alla Fonderia restano documenti e testimoni.\\
\end{tabularx}
\end{overviewbox}

\subsubsection{Situazione iniziale}
La Fonderia rifinisce componenti fabbricati altrove e prepara le spedizioni. La cassa 17, pesata il giorno 13, è rimasta alla Fonderia fino al giorno 16, quando è partita come collo 17-B sul carro di T. Harrow. È il carro seguito da Elias: trasportava il regolatore, vetriatura e minuteria elencati in H05. Qui restano il registro di pesatura H04, la ricevuta H05 e le persone che hanno seguito le operazioni.

Elias aveva rifiutato la richiesta di un intermediario della Chiesa di costruire il regolatore. Una spedizione dello stesso committente aveva attirato la sua attenzione; la differenza fra peso dichiarato e misurato lo aveva spinto a indagare. Non aveva aperto la cassa né riconosciuto il pezzo al suo interno. H01 riporta quel sospetto iniziale. La ricevuta H05 dichiara un regolatore, ma da sola non dimostra che riproduca il suo progetto.

\subsection{Arrivo e ambienti}
\subsubsection{Ingresso}
\begin{readaloud}
I martelli coprono le voci. Accanto al cancello, una squadra spinge una cassa sulla piattaforma della bilancia; un addetto aspetta che l'indice si fermi prima di annotare il peso. Oltre il banco delle consegne, una porta a vetri dà sull'ufficio del caporeparto.

Marta vi chiede di lasciare libero il passaggio dei carrelli. Quando nominate Elias, si toglie un guanto e vi accompagna in ufficio. «Ditemi che cosa cercate.»
\end{readaloud}

\begin{gmbox}
la cassa in lavorazione è una spedizione ordinaria, non la 17. Il rumore rende difficile parlare nel reparto; l'ufficio consente di leggere e conversare. Gli operai stanno lavorando e non conoscono automaticamente il segreto della commessa.
\end{gmbox}

\textbf{Aspetti di scena:} \aspect{Il rumore copre le voci}; \aspect{Carrelli in movimento}; \aspect{Il lavoro degli operai dipende dalle commesse}.

\subsection{Interazioni e ostacoli}
\subsubsection{Parlare con i presenti}
\begin{description}[leftmargin=5mm,style=nextline]
\item[Marta Ghal.]
Conosce le pesature, conserva H04 e sa dove trovare H05. A una domanda su Elias o sulla cassa mostra i documenti in ufficio. Vuole proteggere il lavoro degli operai: chiede di consultare gli originali sul posto e di non accusarla pubblicamente. Non conosce la richiesta dell'intermediario, salvo che i PG le mostrino H01 o gliela raccontino.

Il giorno 16 registrò Elias come presente mentre lui seguiva il carro di T. Harrow, che trasportava la cassa 17-B. Elias voleva verificare che fosse consegnata alla destinazione dichiarata. Marta rischia il posto per la falsa registrazione: lo ammette in privato se i PG spiegano come useranno l'informazione senza esporla. Non vide l'arrivo; conosce il Laboratorio Alchemico come destinazione scritta su H05.
\item[Il contabile della Chiesa.]
Sta verificando le ricevute delle forniture ecclesiastiche. Ha autorizzato la ricodifica 17-B e conosce la descrizione di copertura; ignora lo scopo della ricerca. Se interpellato parla di una procedura riservata. Davanti a H04 e H05 può ammettere che il peso dichiarato era stato riportato dai documenti della commessa, non misurato. Teme di essere accusato di frode; può segnalare le domande ai superiori, ma soltanto quelle di cui viene a conoscenza.
\item[Operai e vetturale.]
Gli addetti possono indicare Marta e il magazzino delle ricevute. T. Harrow, se presente per una consegna o rintracciato tramite Marta, conferma firma e percorso fino al Laboratorio Alchemico. Conosce il contenuto soltanto dalla ricevuta.
\item[Nox, presenza eventuale.]
Inseriscilo se Roland lo ha mandato a verificare il posto di lavoro di Elias o a seguire i PG. Chiede quando Elias sia stato visto e riferisce ciò che apprende. Il suo arrivo deve dipendere dalle ricerche dei Cacciatori, non dalla fine della scena.
\end{description}

\subsubsection{Confrontare pesature e spedizione}
\begin{cluebox}
Con domande pertinenti e consultazione dei documenti, senza tiro, i PG possono stabilire:
\begin{itemize}
\item H04 registra 412 libbre dichiarate e 391 misurate. H01 aggiunge la misura più precisa: 391 e un quarto. La discrepanza è reale, ma non prova un furto.
\item H05 identifica esplicitamente il collo 17-B come la cassa 17. Il peso di 391 libbre e un quarto conferma il confronto.
\item La ricevuta dichiara un regolatore, vetriatura e minuteria e indica il Laboratorio Alchemico come destinazione. Non prova ancora la paternità dell'invenzione.
\item Il riferimento 44-14 in H05 permette di chiedere all'Archivio Ecclesiastico la commessa corrispondente: H06 ne riporta la voce e H06a documenta la spedizione. Il numero individua il documento; le restrizioni di consultazione si gestiscono nel nodo Archivio.
\end{itemize}
\end{cluebox}

\subsubsection{Azioni e complicazioni}
\begin{fatebox}
\textbf{Indagine:} mettere in ordine ricevute e pesature può creare \aspect{Copie e riferimenti pronti al confronto}, utile per contestare la versione del contabile. Leggere il numero di commessa non richiede un tiro.

\textbf{Tecnica:} esaminare le bilance può escludere un guasto e creare \aspect{Pesatura verificata}. Qui non c'è il regolatore da smontare: il suo funzionamento sarà chiarito da H03.

\textbf{Carisma o Empatia:} possono ottenere da Marta discrezione, copie o la confidenza sul turno coperto quando ci sono dubbi e resistenze reali. Un fallimento può attirare l'attenzione del contabile o compromettere la fiducia, senza cancellare le informazioni già ottenute.
\end{fatebox}

\textbf{Compel possibili.} Proponili quando un aspetto pertinente del PG o della scena rende sensata la complicazione, concordandone il costo narrativo con il giocatore.
\begin{itemize}
\item \textbf{Difendere chi lavora.} Un PG sensibile agli abusi accusa il contabile davanti agli operai; il contabile interrompe il colloquio e segnala ufficialmente l'indagine. La ricerca perde discrezione.
\item \textbf{Non lasciare indietro una prova.} Un PG diffidente verso la Chiesa insiste per portar via l'originale affidatogli da Marta. Lei ne risponde personalmente e ritira la propria collaborazione finché non viene restituito o messo al sicuro con il suo accordo.
\item \textbf{Farsi sentire nel reparto.} Se i PG discutono una scoperta sensibile fra i macchinari, il rumore li costringe ad alzare la voce: un'informazione precisa arriva al contabile presente nelle vicinanze. Stabilisci quale, senza attribuirgli l'intera conversazione.
\end{itemize}

\subsection{Handout}
\label{nodo-1-handout}
\phantomsection
\addcontentsline{toc}{subsubsection}{H04 -- Registro di pesatura}
\begin{handoutbox}{H04}{Registro di pesatura}
\textbf{Dove:} Marta lo conserva nell'ufficio della Fonderia.

\textbf{Quando consegnarlo:} Quando i PG chiedono dei controlli di Elias o confrontano i pesi; nessun tiro.

\textbf{Che cosa rivela:} Conferma l'anomalia e l'ordine di non aprire la cassa: 412 libbre dichiarate e 391 misurate.

\textbf{Limiti:} Non contiene da solo la ricodifica o la destinazione della cassa 17.
\end{handoutbox}

\phantomsection
\addcontentsline{toc}{subsubsection}{H05 -- Ricevuta di trasferimento}
\begin{handoutbox}{H05}{Ricevuta di trasferimento}
\textbf{Dove:} Fra le copie carbone del magazzino, sotto il codice 17-B; Marta sa recuperarla.

\textbf{Quando consegnarlo:} Quando i PG chiedono che cosa sia uscito dopo la pesatura e ricostruiscono la spedizione; non occorre trovare la cassa, già partita.

\textbf{Che cosa rivela:} Provenienza, peso, contenuto dichiarato, destinazione e riferimento alla commessa. Il collo 17-B è la cassa 17; il peso accettato è 391 libbre e un quarto.

\textbf{Limiti:} Non indica la sala del Progetto o la responsabile della ricerca e non dimostra da sola la paternità del regolatore.
\end{handoutbox}

\textbf{Se i PG non prendono i documenti:} possono leggerli e annotare i riferimenti. Marta può indicare Laboratorio Alchemico e commessa 44-14 consultando la ricevuta; se questa è stata rimossa, il vetturale conosce almeno la destinazione. Le copie costituiscono prove da mostrare, non l'unico modo di conoscere le piste.

\subsection{Evoluzione della scena}
\begin{clockbox}
Una consultazione discreta non provoca automaticamente una confisca. Se il contabile sente accuse o vede sottrarre documenti, può chiedere un intervento dell'autorità ecclesiastica. Mostra la segnalazione e l'arrivo degli incaricati, lasciando tempo per copiare le carte, trattare o andarsene. Marta salva una copia soltanto se può farlo concretamente. L'eventuale sequestro riguarda i documenti cercati: la Fonderia non interrompe tutta la produzione senza un motivo ulteriore.
\end{clockbox}

\subsection{Collegamenti}
\linkto{Laboratorio Alchemico: destinazione in H05. Archivio Ecclesiastico: originale della commessa 44-14 citato in H05, da confrontare con la scheda H06a. Officina delle Lame: Silas e i vecchi disegni indicati in H01, per verificare l'eventuale uso del progetto di Elias. Rifugio dei Cacciatori: attraverso Nox, se presente.}

\nodehead{11. Nodo 2 -- Studio distrutto e Officina delle Lame}{Nodo biografico; chiarisce la vendita dei disegni e la responsabilità di Silas.}{Come è arrivato il progetto di Elias nelle mani della Chiesa?}

\subsection{A colpo d'occhio}
\begin{overviewbox}
\begin{tabularx}{\linewidth}{>{\bfseries}p{32mm} X}
Presenti & Silas all'Officina delle Lame.\\
Accesso & Studio e officina sono luoghi distinti, poco distanti; si può andare direttamente da Silas.\\
Handout & H03 conservato da Silas, insieme al vecchio esemplare del regolatore.\\
Da ricordare & Silas vendette i disegni sette mesi fa; non fabbricò il pezzo spedito dalla Fonderia.\\
\end{tabularx}
\end{overviewbox}

\subsubsection{Situazione iniziale}
Lo studio incendiato due anni fa e l'Officina delle Lame sono due luoghi distinti, poco distanti. Lo studio mostra le conseguenze della scelta di Elias; l'officina ospita Silas, il frammento H03 e il vecchio esemplare del regolatore. I PG possono recarsi direttamente da Silas: visitare lo studio non è un requisito.

Elias incendiò lo studio dopo aver scoperto che il suo regolatore era stato incorporato in strumenti di tortura. La Chiesa sequestrò frammenti e prototipi, insufficienti a ricostruirlo. Silas aveva conservato una copia di lavoro e la vendette sette mesi fa per pagare debiti, accettando la giustificazione di un impiego ospedaliero. Il giorno 14 confessò la vendita a Elias. Non ha costruito il pezzo spedito dalla Fonderia.

\subsection{Arrivo e ambienti}
\subsubsection{Lo studio bruciato}
\begin{readaloud}
Le pareti dello studio sono annerite, soprattutto attorno alle aperture. Entrando trovate resti di legno bruciato lungo i muri e spazio vuoto dove stavano gli arredi. Non ci sono più attrezzi o disegni da lavoro. Sotto i piedi scricchiolano polvere e piccoli frammenti.
\end{readaloud}

\begin{gmbox}
sono i resti del vecchio incendio, non di un attacco recente. Il materiale recuperabile fu sequestrato dalla Chiesa. Qui non sono previsti nuovi indizi o fogli nascosti; H03 è conservato da Silas. Se i PG cercano, descrivi ciò che resta e chiarisci che la ricerca del progetto deve proseguire attraverso i testimoni.
\end{gmbox}

\textbf{Aspetti di scena:}
\begin{itemize}
\item \aspect{Locali svuotati dopo l'incendio}: c'è poco dietro cui nascondersi; chi entra è facilmente visibile attraverso le aperture.
\item \aspect{Frammenti sotto i piedi}: passi e spostamenti si sentono. Se diventa necessario muoversi senza rumore, i PG possono prima liberare un tratto del pavimento.
\end{itemize}
Usali quando posizione o discrezione contano davvero. La visita non richiede un tiro o una complicazione obbligatoria.

\subsubsection{L'Officina delle Lame}
\begin{readaloud}
Poco distante, dall'officina arrivano colpi di martello e il rumore di una lima. Sul banco ci sono utensili, una morsa e una lama in lavorazione. Silas alza gli occhi, vi indica uno spazio libero accanto al banco e continua a lavorare mentre gli spiegate perché siete venuti.
\end{readaloud}

\begin{gmbox}
quando si nomina l'incendio, Silas depone l'utensile. Quando la conversazione riguarda l'invenzione, può mostrare il vecchio regolatore e H03. Non lasciarli descritti come oggetti misteriosi da individuare con un tiro.
\end{gmbox}

\textbf{Aspetti di scena:}
\begin{itemize}
\item \aspect{Banco attrezzato}: morsa e utensili permettono di esaminare il vecchio regolatore e confrontarlo con H03, con il consenso di Silas.
\item \aspect{Il lavoro copre le voci}: martello e lima ostacolano una conversazione a distanza, ma possono coprire uno scambio a bassa voce vicino al banco. Silas può fermarsi per discutere con chiarezza.
\end{itemize}

\subsection{Interazioni e ostacoli}
\subsubsection{Parlare con Silas}
La sua scheda descrive paure e limiti. Qui vuole aiutare Elias senza esporsi subito al giudizio dei PG. Non conosce il reperto, il trattamento o il ruolo di Celestine.

\begin{description}[leftmargin=5mm,style=nextline]
\item[Che cosa racconta senza resistenze.]
Conosce l'impiego trasfusionale del regolatore e il motivo dell'incendio. Mostra H03 quando gli chiedono del lavoro di Elias. Se vede il marchio amministrativo sui documenti, lo riconosce come quello usato dal compratore ecclesiastico; questo non gli rivela l'identità di Celestine.
\item[Che cosa omette inizialmente.]
Parla del sequestro della Chiesa, ma tace la copia privata che aveva conservato e venduto. Alla domanda su chi fosse il compratore può ammettere di aver avuto rapporti commerciali con lui senza raccontare subito che cosa gli cedette.
\item[Come può arrivare alla confessione.]
Domande precise su come fosse stato ricostruito un progetto incompleto, oppure la spiegazione che la verità può aiutare a trovare Elias, gli danno una ragione per parlare. Empatia può far cogliere la reticenza; non è l'unica via alla confessione. Silas ammette vendita, debiti e giustificazione ospedaliera, precisando che Elias lo sa già dal giorno 14.
\item[Che cosa può offrire.]
H03 documenta il vecchio progetto. Silas può indirizzare i PG a Miriam, a cui presentò Elias per interpretare la runa. Se gli spiegano come intendono usare l'accesso, può consegnare la chiave dell'ingresso di servizio del deposito del Laboratorio Alchemico. Non apre l'ascensore; un buon tiro non ne cambia la funzione.
\end{description}

\textbf{Tre oggetti distinti:} la copia di lavoro fu venduta; H03 è un frammento originale rimasto a Silas, non un altro progetto completo; il vecchio regolatore serve a illustrarne il funzionamento. La ricevuta H05 nomina un regolatore, ma non basta a dimostrare che il pezzo spedito riproduca quello di Elias.

\subsubsection{Azioni e complicazioni}
\begin{fatebox}
\textbf{Tecnica:} il banco, H03 e il vecchio esemplare permettono di creare \aspect{Funzionamento del regolatore compreso}, utile per un successivo confronto. La funzione di base può essere spiegata da Silas senza tiro.

\textbf{Empatia:} cogliere quali domande Silas evita può aiutare a trattare la vendita. Il tiro non obbliga a confessare né sostituisce ciò che il giocatore gli domanda.

\textbf{Carisma:} può sostenere una richiesta concreta di collaborazione, come affidare la chiave o accompagnare i PG al deposito. Prima chiarisci che cosa teme Silas e quali garanzie gli vengono offerte.
\end{fatebox}

\textbf{Compel possibili:} proponili solo se pertinenti a un aspetto del PG e concorda la complicazione; non assegnare ai personaggi convinzioni che non hanno.
\begin{itemize}
\item \textbf{La lealtà a Elias rende difficile ascoltare.} Un PG legato a Elias reagisce alla confessione accusando Silas di averlo tradito. Silas ritira l'offerta della chiave e di accompagnarli; per recuperare la collaborazione occorre ricucire il rapporto. La confessione già ascoltata resta acquisita.
\item \textbf{Conservare la prova a ogni costo.} Un PG convinto che i documenti finiscano sempre nelle mani sbagliate decide di trattenere H03 senza l'accordo di Silas. Silas se ne accorge, interrompe il colloquio e chiede la restituzione. Conservare il frammento compromette la sua fiducia e l'aiuto volontario; restituirlo permette di tentare una riconciliazione.
\end{itemize}

\subsection{Handout}
\label{nodo-2-handout}
\phantomsection
\addcontentsline{toc}{subsubsection}{H03 -- Progetto del regolatore}
\begin{handoutbox}{H03}{Progetto del regolatore}
\textbf{Dove:} Il frammento originale è conservato da Silas all'Officina delle Lame.

\textbf{Quando consegnarlo:} Silas lo mostra quando i PG chiedono dell'invenzione, discutono H01 o portano H05. Non devono avere recuperato un componente alla Fonderia: la cassa è già partita.

\textbf{Che cosa rivela:} L'origine trasfusionale e la funzione del regolatore di Elias. La confessione di Silas spiega separatamente come la Chiesa abbia ottenuto i disegni. Il confronto con il pezzo attuale avverrà nella sala del Progetto.

\textbf{Limiti:} Non descrive la macchina costruita dalla cerchia, il trattamento o lo scopo di Celestine.
\end{handoutbox}

\textbf{Se H03 non viene preso:} Silas può spiegare il progetto e raccontare la vendita. I PG rinunciano alla prova materiale da confrontare in seguito, non alle informazioni. La consegna della chiave è una decisione distinta e non dipende dal possesso di H03.

\subsection{Evoluzione della scena}
\begin{clockbox}
L'arresto di Silas non scatta alla fine della scena. Può diventare una reazione della Chiesa solo se questa viene a conoscenza della sua collaborazione o delle indagini sulla copia venduta. Stabilisci chi ha riferito che cosa e mostra l'intervento, lasciando ai PG possibilità di agire. Se avviene l'arresto, testimoni possono indicare la Caserma; Silas resta raggiungibile. Una visita discreta non provoca automaticamente conseguenze.
\end{clockbox}

\subsection{Collegamenti}
\linkto{Fonderia per confrontare le spedizioni; Archivio tramite il marchio delle commesse; Biblioteca tramite la presentazione a Miriam; Laboratorio Alchemico tramite l'ingresso di servizio e la relativa chiave.}

\nodehead{12. Nodo 3 -- Archivio Ecclesiastico}{Consultazioni, commesse e disposizioni su Elias.}{Che cosa chiedono i PG, e fin dove Ilyne è disposta ad aiutarli?}

\subsection{A colpo d'occhio}
\begin{overviewbox}
\begin{tabularx}{\linewidth}{>{\bfseries}p{32mm} X}
Presenti & Ilyne, copisti, visitatori e due custodi. Aldren solo se contattato o avvisato.\\
Accesso & Richiesta circostanziata: H05 oppure dati di cassa, data e pesi da H01/H04. Consultazione al tavolo, originali custoditi. Il contenuto di H07 richiede autorizzazione di Aldren, collaborazione di Ilyne o accesso clandestino concreto.\\
Handout & H06: riepilogo; H06a: spedizione 44-14; H06b: disegno 44-12, su domanda pertinente; H07: disposizione riservata.\\
Da ricordare & L'Archivio conserva e verifica le commesse delle strutture della Chiesa. Ilyne segnala automaticamente l'esistenza di H07 quando i PG dichiarano di cercare Elias; ottenerlo non è necessario per proseguire.\\
\end{tabularx}
\end{overviewbox}

\textbf{Uso al tavolo:} passa alla situazione corrispondente a ciò che fanno i PG; i blocchi non sono fasi obbligatorie. Gli aspetti descrivono opportunità e limiti visibili. Proponi i compel soltanto quando pertinenti, concordando la complicazione: non sono eventi automatici né condizioni per ricevere gli indizi ordinari.

\subsection{Arrivo e ambienti}
\subsubsection{Ingresso}
\begin{readaloud}
Davanti al banco aspettano alcuni visitatori con ricevute ripiegate in mano. Ilyne controlla un riferimento, lo ripete al copista e indica il tavolo dove attendere. Le voci rimangono basse, interrotte dal rumore delle pagine e delle sedie spostate.

Dietro il banco, una porta conduce agli scaffali. Quando Ilyne va a recuperare un registro, il custode lascia passare lei e ferma con un gesto il visitatore che prova a seguirla. Dal banco sentite le richieste di chi vi precede. Quando arriva il vostro turno, Ilyne posa la penna: «Quale pratica volete verificare?»
\end{readaloud}

\textbf{Spazi visibili:} il banco riceve le richieste; al tavolo si leggono i documenti sotto supervisione; oltre la porta c'è il deposito riservato al personale. Ilyne vi conserva H07 fra le disposizioni interne separate dai registri ordinari. I custodi controllano il passaggio e impediscono di portare via originali.

\textbf{Al banco:} \aspect{Richieste udibili da chi aspetta}. I PG possono ascoltare uno scambio o rendere pubblica una contestazione; anche le loro parole possono essere sentite. Ascoltare una frase pronunciata accanto a loro non richiede un tiro.

\subsection{Interazioni e ostacoli}
\subsubsection{Verificare le forniture}
\textbf{Ilyne risponde:} chiede il riferimento della spedizione. H05 indica 44-14; se manca, Fonderia, cassa 17, data e pesi di H01/H04 bastano a rintracciarla. Non occorre fingersi funzionari o superare una prova di Carisma. Con il solo nome di Elias chiede quale fornitura intendano verificare; passa anche al blocco successivo.

\textbf{Rimando agli handout:} H06 per il riepilogo, H06a per la spedizione, H06b per il disegno della vasca. Le schede nella sezione Handout distinguono le richieste che li rendono consultabili.

\begin{fatebox}
\textbf{Aspetto al tavolo:} \aspect{Consultazione sotto gli occhi di Ilyne}. La sua presenza rende difficile sottrarre o alterare fogli, ma permette di ottenere una conferma autorevole.

\textbf{Azione possibile:} confrontando H04 e H06a con Ilyne, i PG possono creare con Indagine \aspect{Discrepanza verificata dall'archivista}, utile per sostenere una richiesta ad Aldren. La conferma riguarda i pesi scritti, non la prova di un furto o della colpa di Celestine. Leggere i numeri resta gratuito e automatico; se il confronto non crea un vantaggio, le informazioni già lette restano note.
\end{fatebox}

\subsubsection{Chiedere di Elias}
\textbf{Ilyne sa:} il giorno 18 Elias si presentò come verificatore delle pesature della Fonderia, ruolo riportato in H04. Consultò le commesse e il disegno della vasca, poi copiò il segno nei suoi appunti.

\textbf{Quando dichiarano di cercarlo, dice automaticamente:}
\begin{readaloud}
«Ward è venuto qui a verificare queste commesse. Dopo la sua visita abbiamo ricevuto una disposizione che lo riguarda. Posso mostrarvi i documenti delle forniture; per quell'ordine occorre l'autorizzazione di Padre Aldren.»
\end{readaloud}

\textbf{Che cosa ottengono:} l'esistenza di H07, il fatto che sia successivo alla visita e il nome a cui chiedere accesso. Ilyne non rivela ancora l'ordine di trattenere Elias e recuperare le sue carte. Se chiedono soltanto delle merci, non introduce spontaneamente H07. L'ordine riguarda Elias, non impone di fermare chiunque lo cerchi.

\begin{fatebox}
\textbf{Compel al banco:} se discutono apertamente della ricerca, \aspect{Richieste udibili da chi aspetta} può far perdere discrezione. Proponi che un visitatore senta una frase precisa e la riferisca ad Aldren; concorda quale informazione viene divulgata e quando può raggiungerlo. Non conosce il resto dell'indagine. Se i PG hanno ottenuto un colloquio privato, questo appiglio non si applica.
\end{fatebox}

\subsubsection{Ottenere accesso all'ordine riservato}
\paragraph{Chiedono l'autorizzazione di Aldren}
\textbf{Risposta:} Aldren vuole controllare l'indagine e può chiedere aggiornamenti o copie delle loro carte in cambio. Non è già presente senza motivo: i PG devono contattarlo oppure attendere che venga avvisato.

\textbf{Esito:} Ilyne mostra H07 se il permesso riguarda le disposizioni su Elias. Un'autorizzazione a verificare una spedizione non basta.

\textbf{Prova eventuale:} Carisma per negoziare condizioni precise quando Aldren oppone una resistenza. Un rifiuto lascia aperte le altre piste; non sottrae ai PG documenti mai consegnati a lui.

\paragraph{Chiedono a Ilyne di aiutarli}
\textbf{Risposta:} cercare uno scomparso e voler sapere se sia stato trattenuto è una motivazione pertinente. Ilyne chiede come useranno l'ordine e come eviteranno di esporla. Non pretende che abbiano già dimostrato l'esperimento.

\textbf{Esito:} se sceglie di collaborare, mostra H07 in privato e può permetterne una copia. Una promessa generica di rispettare la Chiesa non è sufficiente.

\begin{fatebox}
\textbf{Prova eventuale:} Empatia per comprendere il timore di perdere il posto e la possibilità di vigilare sui registri; Carisma per sostenere una proposta concreta. Un tiro non obbliga alla fiducia. Se la richiesta non la convince, può rimandare ad Aldren.

\textbf{Compel sulla riservatezza:} con un aspetto pertinente di lealtà o protezione, un PG può impegnarsi a non esporre Ilyne come fonte. La complicazione è che esibire pubblicamente H07 potrebbe rivelare il suo aiuto: concorda che cosa comporta la promessa. Non rendere il compel l'unico modo di ottenere la sua collaborazione.
\end{fatebox}

\paragraph{Tentano di raggiungere l'ordine senza permesso}
\textbf{Che cosa vedono:} \aspect{Originali oltre la porta controllata}. Ilyne porta e riporta i fogli dal deposito; una distrazione al banco non elimina il custode della porta.

\textbf{Prima del tiro:} devono individuare dove sono conservate le disposizioni e quale controllo eludere. Conoscere l'esistenza di H07 non ne rivela il posto esatto. Se osservano un movimento effettivo di documenti, stabilisci che cosa possono seguirne; non far comparire l'ordine solo per offrire un'occasione.

\textbf{Azioni possibili:} Inganno per un incarico o una motivazione falsa esplicitamente dichiarati; Furtività per raggiungere o copiare un documento individuato. Una distrazione può creare un vantaggio riferito all'attenzione di una persona concreta, non concedere l'intero archivio.

\textbf{Se scoperti:} Ilyne sospende la consultazione e i custodi fermano il passaggio o la sottrazione. Un'informazione mai letta non viene concessa automaticamente; le copie già ottenute restano un fatto della scena.

\subsection{Handout}
\label{nodo-3-handout}
\phantomsection
\addcontentsline{toc}{subsubsection}{H06 -- Registro delle commesse}
\begin{handoutbox}{H06}{Registro delle commesse}
\textbf{Dove:} Nei registri custoditi dall'Archivio; Ilyne lo porta al tavolo di consultazione.

\textbf{Quando consegnarlo:} Quando viene avviata una verifica con un riferimento circostanziato, come descritto in Interazioni e ostacoli.

\textbf{Che cosa rivela:} Materiali, destinatari e autorizzazione comune delle commesse.

\textbf{Limiti:} Il marchio dell'ufficio riservato non vieta la consultazione della fornitura.
\end{handoutbox}

\phantomsection
\addcontentsline{toc}{subsubsection}{H06a -- Spedizione 44-14}
\begin{handoutbox}{H06a}{Spedizione 44-14}
\textbf{Dove:} Nella documentazione della commessa conservata all'Archivio.

\textbf{Quando consegnarlo:} Quando i PG chiedono della commessa, del contenuto della cassa o del percorso.

\textbf{Che cosa rivela:} Il collegamento 17/17-B, peso reale di 391 libbre e un quarto, contenuto, vetturale, data e consegna al deposito del Laboratorio Alchemico. Recupera la pista di H05 se la ricevuta manca e la conferma se è disponibile.

\textbf{Limiti:} Il peso riguarda l'intero collo.
\end{handoutbox}

\phantomsection
\addcontentsline{toc}{subsubsection}{H06b -- Disegno 44-12}
\begin{handoutbox}{H06b}{Disegno 44-12}
\textbf{Dove:} Come foglio separato allegato alla commessa della vasca.

\textbf{Quando consegnarlo:} Quando i PG chiedono della vasca, degli allegati o di ciò che consultò Elias. Non darlo automaticamente insieme al registro.

\textbf{Che cosa rivela:} La vasca e la runa da incidere. Ilyne ricorda che Elias copiò la runa.

\textbf{Limiti:} Ilyne non interpreta il segno.
\end{handoutbox}

\textbf{Consultazione dei tre documenti ordinari:} i PG possono leggere e prendere appunti; chiedono a Ilyne per copiare i fogli. Gli originali restano all'Archivio. Le tre carte possono emergere nella stessa visita, senza tre prove successive. Non rivelano Celestine, il trattamento o la destinazione sotterranea del regolatore.

\phantomsection
\addcontentsline{toc}{subsubsection}{H07 -- Disposizione interna riservata}
\begin{handoutbox}{H07}{Disposizione interna riservata}
\textbf{Dove:} Fra le disposizioni interne nel deposito riservato al personale, separatamente dai registri ordinari.

\textbf{Quando consegnarlo:} Soltanto quando Ilyne lo mostra su autorizzazione pertinente di Aldren, sceglie di collaborare in privato, oppure i PG raggiungono effettivamente il documento. Le condizioni sono descritte in Interazioni e ostacoli. Ilyne può permetterne una copia se collabora.

\textbf{Che cosa rivela:} L'ordine di trattenere Elias con cortesia, avvisare Aldren e recuperare le sue carte; richiama trasferimenti attraverso passaggi di servizio.

\textbf{Limiti:} La sola notizia dell'esistenza dell'ordine non ne rivela il contenuto. Non identifica Celestine né dimostra che Aldren diriga il Progetto. H07 non è necessario per proseguire.
\end{handoutbox}

\subsection{Evoluzione della scena}
\begin{description}[leftmargin=5mm,style=nextline]
\item[I PG contestano pubblicamente una restrizione.]
Ilyne spiega il limite della consultazione. Se insistono fino a interrompere il lavoro, può sospenderla e chiedere l'intervento di Aldren. \textbf{Compel possibile:} un aspetto di orgoglio o insofferenza verso l'autorità spinge il PG a rendere pubblica la contestazione, con quella conseguenza concordata. Non scatta un arresto automatico.
\item[Arriva un ordine di ritiro.]
Usa \aspect{Ogni ritiro deve essere registrato}. Stabilisci chi ha disposto la rimozione, quali carte riguarda e dove vengono portate; Ilyne lo annota. I PG possono chiedere che la procedura venga rispettata e invocare l'aspetto, se pertinente a una contesa sull'occultamento del ritiro. Le registrazioni dovute non costano punti Fate. Nessuna pagina viene strappata alla semplice fine della scena.
\item[Aldren interviene.]
Deve essere stato contattato o informato. Sa soltanto ciò che gli è stato riferito; può offrire il proprio aiuto in cambio degli appunti. Registra quali carte i PG gli affidano.
\item[I PG rinunciano a H07 o lasciano l'Archivio.]
H06a, come H05, conduce al Laboratorio Alchemico; H06b offre la runa da confrontare con Miriam. H07 aggiunge informazioni sulle azioni contro Elias, non è una tappa obbligatoria. Se manca una richiesta verificabile, possono tornare con i riferimenti della Fonderia.
\end{description}

\subsection{Collegamenti}
\linkto{Laboratorio Alchemico tramite H06a o H05; Ospedale Vecchio tramite il registro; Biblioteca Proibita tramite H06b; Padre Aldren per le disposizioni su Elias.}

\nodehead{13. Nodo 4 -- Ospedale Vecchio}{Visita di Elias, testimonianze sui trattamenti e accesso all'ala inferiore.}{Che cosa è successo ai pazienti e chi può aiutare i PG a verificarlo?}

\subsection{A colpo d'occhio}
\begin{overviewbox}
\begin{tabularx}{\linewidth}{>{\bfseries}p{32mm} X}
Presenti & Oren Sile, Maud Kerren, personale e pazienti. Aldren solo se la sua presenza deriva dagli eventi di gioco.\\
Accesso & Reparto pubblico e colloqui con il personale; la scala interna per l'ala inferiore è oltre una porta sorvegliata. Oren e Maud conoscono l'ingresso, ma non possono autorizzare estranei.\\
Handout & H08 presso Oren; H09 fra le sue carte personali nell'ufficio. Maud conserva separatamente la lista dei nomi.\\
Da ricordare & Oren visitò Elias il giorno 19: era malato ma lucido; H08 non certifica le condizioni attuali. I trattamenti danno benefici temporanei e ricadute; sedute più lunghe e frequenti hanno portato dal consenso alle imposizioni. Sei pazienti sono stati trasferiti contro la loro volontà il giorno 25; due nomi sono cancellati dai registri pubblici.\\
\end{tabularx}
\end{overviewbox}

\textbf{Uso al tavolo:} scegli il blocco corrispondente alle domande e alle azioni dei PG; non sono passaggi obbligatori. Le informazioni emergono dalle persone e dai documenti indicati, non automaticamente entrando nell'edificio. Proponi i compel soltanto quando pertinenti a un aspetto, concordando la complicazione.

\subsection{Arrivo e ambienti}
\subsubsection{Il reparto pubblico}
\begin{readaloud}
Il reparto è affollato e continua a funzionare. Fra i letti il personale visita i malati e porta avanti le cure; per parlare con qualcuno bisogna prima attirarne l'attenzione e lasciargli finire ciò che sta facendo. Oren lavora qui, fra il reparto e il suo ufficio. Maud chiama i pazienti per nome.
\end{readaloud}

\textbf{Come iniziare:} chi cerca Oren può chiederne al personale, spiegando di cercare Elias Ward. Maud può indicare dove trovarlo. Questo permette un colloquio, non autorizza a consultare cartelle altrui o a entrare nell'ufficio senza essere ricevuti.

\textbf{Aspetto:} \aspect{Reparto sovraffollato}. Il movimento fra i letti può aiutare a passare inosservati, ma rende difficile una conversazione riservata. Per ottenere privacy i PG devono scegliere dove e quando parlare: l'ufficio di Oren è un'opzione se lui li riceve. Nessun tiro è necessario per chiederglielo.

\subsection{Interazioni e ostacoli}
\subsubsection{Chiedere della visita di Elias}
\textbf{Oren risponde:} conosce Elias da precedenti collaborazioni mediche e lo ha visitato il giorno 19. Descrive insonnia, irritabilità, sensibilità agli odori e primi segni bestiali, insieme a linguaggio coerente e memoria integra. Elias comprendeva il proprio stato, temeva un peggioramento e aveva rifiutato altra emoterapia. Chiese anche delle apparecchiature ricevute dall'ala riservata.



\textbf{Rimando agli handout:} H08 documenta la visita; consegna e limiti sono raccolti nella sezione Handout.

\begin{gmbox}
la malattia di Elias precede il suo ingresso nei sotterranei. H08 documenta la lucidità al giorno 19, non esclude un peggioramento successivo e non dimostra da solo la verità dei suoi sospetti. Oren non sa dove si trovi ora.
\end{gmbox}

\textbf{Medicina o Emoterapia:} possono servire per interpretare i sintomi e discutere la valutazione con Oren. La data, le osservazioni scritte e il fatto che il paziente fosse lucido sono leggibili senza tiro. Un eventuale vantaggio riguarda la valutazione documentata, non una diagnosi a distanza delle condizioni attuali.

\subsubsection{Chiedere dei pazienti trasferiti}
\textbf{Maud risponde:} conosce i sei pazienti e li assisteva prima del trasferimento. Il giorno 25 furono accompagnati nell'ala inferiore malgrado avessero rifiutato altre sedute. Confrontando i registri, ha scoperto che due dei loro nomi non compaiono più. Non sono sei persone di cui nessuno conosce la destinazione: è la documentazione pubblica a nascondere parte di ciò che è accaduto.

\textbf{La lista di Maud:} ha conservato una copia dei nomi invece di distruggerla come le era stato ordinato. La mostra a chi cerca i malati e spiega come intende usare l'informazione. Non occorre convincerla a desiderare il loro soccorso; teme di perdere il posto e di non poterli più aiutare. La lista è distinta da H09 e non prova che i due cancellati siano morti.

\begin{fatebox}
\textbf{Appiglio:} i nomi conservati da Maud permettono un confronto preciso con i registri. Se i PG ottengono accesso a quei registri, Indagine può servire a creare un vantaggio basato sul confronto; non è necessaria per capire il racconto dell'infermiera.

\textbf{Compel possibile:} se un aspetto del PG esprime insofferenza per gli abusi, proponi che contesti ad alta voce il trasferimento. La complicazione è rendere udibile agli altri la conversazione che Maud voleva mantenere privata. Non implica automaticamente una denuncia o l'arrivo di Aldren.
\end{fatebox}

\subsubsection{Capire gli effetti dei trattamenti}
\textbf{Fatti osservati:} Oren e Maud hanno visto ritorni di lucidità seguiti da ricadute nello stato precedente. Inizialmente i pazienti erano informati e potevano rifiutare. Le sedute sono poi diventate più lunghe e ravvicinate; le proteste sono state ignorate. Maud può riferire ciò che ha visto e sentito dai malati, senza spiegare il funzionamento della macchina.

\textbf{Limite delle conoscenze di Oren:} ha seguito gli effetti clinici fin dalla fase volontaria, ma ha visto la vasca nella sala del Progetto soltanto il giorno 25. Può descrivere il sangue che passa dal paziente alla macchina e al sistema della vasca, per poi tornare al paziente, e la presenza del reperto vivo. Non sa da dove provenga il reperto, non conosce la finalità ultraterrena e non sa spiegare perché il beneficio svanisca. Vedere l'impianto non lo rende capace di manovrarlo o di indicare un'interruzione sicura della seduta.

\textbf{Documentazione clinica:} Oren può confrontare i referti con i PG. Leggere le date e riconoscere sedute sempre più frequenti non richiede Medicina; la competenza serve a interpretare le osservazioni cliniche. I referti e le testimonianze documentano effetti temporanei e coercizione, non dimostrano il contatto con un'entità. La ragione per cui Celestine insiste va ricostruita dalle altre fonti.

\subsubsection{Affrontare le responsabilità di Oren}
\textbf{Cosa nasconde:} il giorno 21 informò Aldren della visita di Elias, sperando in una verifica delle prove e nella sua protezione. Il giorno 25 eseguì l'ordine di accompagnare i sei pazienti nell'ala inferiore contro la loro volontà. Teme incriminazione, rappresaglie e il disprezzo di Eleanor.

\textbf{Come parlarne:} inizialmente presenta la denuncia come una richiesta di consulto. H07, se disponibile, permette di contestarne le conseguenze. Una conversazione privata e una proposta concreta per aiutare i pazienti senza abbandonare i testimoni possono indurlo ad assumersi la responsabilità. H07 non è un requisito obbligatorio e i PG non devono assolverlo.

\textbf{Empatia:} può aiutare a cogliere il timore quando si nomina Aldren o l'esitazione sul trasferimento. Non rivela automaticamente fatti nascosti e non obbliga Oren a confessare; l'eventuale confronto si gioca sulla base di ciò che viene detto e mostrato.

\textbf{Rimando agli handout:} H09 può sostenere la confessione di Oren; la scheda precisa custodia e condizioni di consultazione.

\subsubsection{Scendere nell'ala inferiore}
\begin{readaloud}
La porta che conduce alla scala dell'ala inferiore è chiusa e sorvegliata. Alla richiesta di passare, viene risposto che il settore è chiuso per lavori. Le tubazioni nuove proseguono verso quella parte dell'edificio. Da qui non si vedono né le stanze dei pazienti trasferiti né la sala del Progetto.
\end{readaloud}

\textbf{Percorso noto:} Oren e Maud possono indicare la porta senza una prova di Indagine. La scala conduce alle stanze dei pazienti nell'ala inferiore; sullo stesso livello, il corridoio porta alla sala del Progetto. Non c'è un ulteriore laboratorio nascosto a un piano ancora più basso. Le tubazioni sono osservabili, ma da sole non provano un abuso.

\textbf{Ostacolo concreto:} la sorveglianza controlla chi scende. L'accesso al reparto pubblico o il consenso a parlare con Oren non autorizzano a oltrepassare la porta. Occorre un'autorizzazione valida per l'ala inferiore oppure un tentativo concreto di aggirare o affrontare il controllo. Oren e Maud possono spiegare il percorso, non concedere il permesso; un'eventuale scorta del personale non sostituisce l'autorizzazione.

\begin{fatebox}
\textbf{Aspetto:} \aspect{Accesso sorvegliato all'ala inferiore}. Rende il passaggio un ostacolo visibile e permette ai PG di osservare il controllo prima di agire. Se hanno un'autorizzazione pertinente, la mostrano; se cercano un pretesto o una distrazione, stabilisci chi fa cosa e quale opposizione incontra. Una prova riguarda quel tentativo, non l'intera esplorazione dell'ala.

\textbf{Creare un vantaggio:} osservare il controllo può preparare un tentativo successivo. Stabilisci ciò che i PG possono effettivamente vedere: un successo non crea una chiave, un passaggio segreto o una pausa delle guardie.
\end{fatebox}

\subsection{Handout}
\label{nodo-4-handout}
\phantomsection
\addcontentsline{toc}{subsubsection}{H08 -- Referto di visita non protocollato}
\begin{handoutbox}{H08}{Referto di visita non protocollato}
\textbf{Dove:} Oren conserva il referto della visita del giorno 19.

\textbf{Quando consegnarlo:} Quando il colloquio affronta la visita di Elias. Oren lo mostra a chi lo sta cercando; non richiede una prova di Carisma o la ricerca di un archivio medico.

\textbf{Che cosa rivela:} Sintomi e lucidità di Elias al giorno 19, prima dell'ingresso nei sotterranei.

\textbf{Limiti:} Non certifica le condizioni attuali di Elias, non dimostra da solo la verità dei suoi sospetti e non contiene analisi del sangue o istruzioni sul Progetto.
\end{handoutbox}

\phantomsection
\addcontentsline{toc}{subsubsection}{H09 -- Nota mai spedita}
\begin{handoutbox}{H09}{Nota mai spedita}
\textbf{Dove:} Fra le carte personali nell'ufficio di Oren. La conserva perché non sa a chi affidarla senza esporsi o esporre altri.

\textbf{Quando consegnarlo:} Quando Oren decide di raccontare la propria responsabilità e cercare aiuto. Per prenderla senza di lui occorre stabilire come i PG entrano nell'ufficio e raggiungono le carte; non basta dichiarare una ricerca nell'Ospedale. Maud può consegnarla soltanto se Oren gliel'ha affidata, per esempio prima di una convocazione di Aldren: registra il passaggio se avviene.

\textbf{Che cosa rivela:} Benefici e ricadute, sedute imposte, denuncia del giorno 21 e trasferimento del 25. Indica l'ala inferiore e descrive il collegamento fra il corridoio e la sala del Progetto.

\textbf{Limiti:} Non spiega l'origine del reperto, il criterio di scelta dei pazienti o l'ipotesi di Celestine sul contatto.
\end{handoutbox}

\textbf{Altro documento:} la lista conservata da Maud resta distinta da H09; uso e condizioni sono descritti nel colloquio sui pazienti trasferiti.

\textbf{Informazioni senza handout:} Oren può raccontare la visita senza consegnare H08 e ammettere i fatti di H09 senza cedere la lettera. Maud può riferire nomi e trasferimento. Mancano le prove scritte da mostrare ad altri, non necessariamente le informazioni. Le confidenze di Oren richiedono comunque le condizioni descritte sopra.

\subsection{Evoluzione della scena}
\begin{clockbox}
Non far scattare una nuova chiusura solo perché i PG arrivano: la porta è già chiusa e sorvegliata. Se dagli eventi nasce un ordine che impedisce a Maud di avvicinarsi ai sei pazienti, rendi chiaro chi lo impartisce e perché; lei può cercare aiuto e indicare l'ingresso. Se Oren prepara la fuga, Maud può segnalarlo. Non introdurre altri trasferimenti o una trasformazione improvvisa per obbligare i PG a intervenire.
\end{clockbox}

\subsection{Collegamenti}
\linkto{Laboratorio Alchemico per le forniture e l'accesso dal deposito; Archivio Ecclesiastico per confrontare le prove con le commesse e H07; Rifugio dei Cacciatori per discutere la valutazione di Elias; Biblioteca Proibita per confrontare quanto osservato con la lettura della runa. L'accesso interno all'ala inferiore può condurre direttamente al nodo finale se i PG superano il controllo.}

\nodehead{14. Nodo 5 -- Rifugio dei Cacciatori}{Ricerca di Elias e accordi per un recupero sorvegliato.}{Che cosa sanno i Cacciatori e quali garanzie chiedono per aiutare i PG?}

\subsection{A colpo d'occhio}
\begin{overviewbox}
\begin{tabularx}{\linewidth}{>{\bfseries}p{32mm} X}
Presenti & Roland, Nox e cacciatori feriti, salvo incarichi o spostamenti già avvenuti in gioco.\\
Accesso & Colloquio con Roland e consultazione dell'avviso sulla bacheca. Un eventuale recupero richiede un accordo su tempi, aiuto e garanzie.\\
Handout & H10 sulla bacheca o mostrato da Roland; non è una diagnosi né un ordine di esecuzione preventiva.\\
Da ricordare & La ricerca è iniziata il giorno 27 dopo la segnalazione ospedaliera sull'intrusione del 26. Il sorvegliante riferisce che Elias lo ha colpito al braccio per liberarsi e lo ha risparmiato. I Cacciatori non conoscono esperimenti, sala del Progetto o nascondiglio.\\
Aiuto possibile & Tempo concordato, scorta, contenimento e protezione delle vie di evacuazione.\\
\end{tabularx}
\end{overviewbox}

\textbf{Uso al tavolo:} passa al blocco corrispondente alle domande dei PG. Le informazioni ordinarie non richiedono una sequenza di prove. Un tiro risolve un disaccordo o un'azione concreta; i compel sono proposte legate agli aspetti, non eventi obbligatori.

\subsection{Arrivo e ambienti}
\subsubsection{Ingresso}
\begin{readaloud}
Nella sala alcuni Cacciatori puliscono le armi; altri, feriti, riposano sulle sedute. Sulla bacheca è affisso un avviso con il nome di Elias Ward. Una mappa è aperta sul tavolo dove Roland esamina le informazioni sulla ricerca. Nox guarda verso di lei prima di intervenire.

Quando indicate l'avviso o chiedete del ricercato, Roland posa le carte: «Che cosa sapete di Elias Ward?»
\end{readaloud}

\textbf{Mappa sul tavolo:} è uno strumento delle pattuglie, non un indizio da decifrare o un handout. Non indica una pista di Elias nelle fogne, l'ingresso dello scarico o un percorso verso il Progetto. Non occorre superare una prova per ottenere da Roland le informazioni descritte sotto.

\textbf{Aspetto:} \aspect{Una ricerca già in corso}. I PG possono proporre di coordinarsi con le pattuglie, ma devono concordare che cosa condividere e chi farà che cosa. I Cacciatori conoscono soltanto luoghi e accessi effettivamente individuati; non presidiano automaticamente ogni uscita.

\subsection{Interazioni e ostacoli}
\subsubsection{Chiedere perché cercano Elias}
\textbf{Roland risponde:} il giorno 27 l'Ospedale ha segnalato un intruso malato e armato. L'episodio era avvenuto il 26 nel settore del deposito interno al Laboratorio Alchemico. Il sorvegliante ferito le ha raccontato che Elias lo colpì al braccio per liberarsi e fuggì senza finirlo, pur avendone la possibilità. Roland non ha visto Elias e non dichiara che sia già una bestia; teme il rischio di una trasformazione.

\begin{gmbox}
Celestine ha organizzato la segnalazione omettendo la ricerca. Roland conosce la provenienza ospedaliera, non questa regia. Non può indicare Celestine come responsabile o spiegare gli esperimenti senza ricevere nuove prove.
\end{gmbox}



\textbf{Il sorvegliante:} è la guardia ferita al deposito, non uno dei Cacciatori feriti nella sala. I PG possono chiederne notizie a Roland e rintracciarlo al posto di cura dell'Ospedale. La sua testimonianza conferma il gesto con cui Elias lo ha risparmiato; non garantisce il comportamento futuro. Perse di vista Elias dopo lo scontro e non lo vide prendere l'ascensore.

\textbf{Rimando agli handout:} H10 riporta l'ordine; la scheda nella sezione Handout ne precisa la consegna.

\subsubsection{Chiedere dove cercarlo}
\textbf{Informazione disponibile:} il punto certo è l'intrusione al deposito del Laboratorio Alchemico il giorno 26. Roland può riferire il luogo dell'episodio e indirizzare al testimone; non conosce il successivo percorso sotterraneo di Elias. Un'indagine sul deposito rimanda al nodo 7, senza consegnare qui un accesso alla sala del Progetto.

\textbf{Limite della pista:} non attribuire alle pattuglie ritrovamenti di trappole, ostacoli o tracce nelle fogne. Se durante il gioco ottengono una nuova pista, registra chi l'ha trovata e come l'ha comunicata. La ricerca dei Cacciatori non risolve fuori scena gli accessi descritti in H14.

\subsubsection{Discutere la pericolosità di Elias}
\textbf{Distinguere le prove:} H08 documenta lucidità al giorno 19. Il racconto del sorvegliante riguarda un gesto di autocontrollo il giorno 26. Sono ragioni per tentare un recupero e ascoltare Elias, non garanzie sulle sue condizioni attuali. Roland può riconoscerne il valore senza rinunciare alle precauzioni.

\textbf{Medicina o Emoterapia:} aiutano a interpretare sintomi e referti, ma non producono una diagnosi attuale senza esaminare Elias. La capacità di usare strumenti dimostra competenza, non basta a provare autocontrollo. Un incontro nel quale risponde, si ferma a richiesta e coopera offre elementi nuovi.

\textbf{Empatia:} può far cogliere che Roland teme un recupero senza misure di sicurezza. In passato esitò davanti a una persona già trasformata che sembrava riconoscerla; l'aggressione successiva costò cinque vite. Può raccontarlo in privato a chi prende sul serio il problema, ma non è una confessione ottenuta automaticamente con un tiro né una condizione per negoziare.

\begin{fatebox}
\textbf{Aspetto:} \aspect{Contenerlo, se possibile}. I PG possono richiamare Roland e Nox al loro stesso ordine e sostenere un tentativo di dialogo protetto. Leggere l'ordine non richiede punti Fate; un'invocazione sostiene una prova pertinente quando esiste opposizione.

\textbf{Creare un vantaggio:} confrontare H08 e la testimonianza con Roland può sostenere \aspect{Prove di autocontrollo da verificare sul posto}. Il vantaggio aiuta a negoziare un tentativo di recupero; non stabilisce in anticipo come reagirà Elias.
\end{fatebox}

\subsubsection{Concordare un recupero}
\textbf{Le domande di Roland:}
\begin{itemize}
\item Quanto tempo chiedete, e dove ci ritroviamo per un aggiornamento?
\item Come intendete avvicinarlo e parlargli?
\item Chi proteggerà i presenti e terrà libera l'uscita se perde il controllo?
\end{itemize}

\textbf{L'accordo:} registra termine, luogo di incontro, informazioni da condividere e aiuto richiesto. Roland può concedere tempo, accompagnare i PG o proteggere le vie di evacuazione. Non occorre un tiro per porre la richiesta: Carisma può risolvere un disaccordo specifico sulla durata, sulla scorta o sulle garanzie. Un fallimento riguarda quel punto della trattativa, non cancella le informazioni già disponibili.

\textbf{Nox:} vuole dimostrarsi affidabile. Un incarico preciso, come tenere libera un'uscita, riconosciuto da Roland gli dà un compito concreto. Conosce H10 e ciò che ha appreso di persona o dagli ordini ricevuti; non sa automaticamente ciò che sanno la capitana o i PG.

\textbf{Limiti dell'intervento:} Roland rispetta l'accordo e valuta i fatti osservati. Aspetto alterato, rabbia o spasmi non bastano per colpire. Se Elias aggredisce senza riuscire a fermarsi e non può essere contenuto, interviene per proteggere i presenti; quando la scena lo consente, annuncia il rischio e lascia spazio a una reazione. L'alleanza non le attribuisce automaticamente il diritto di decidere da sola il destino di Elias.

\begin{fatebox}
\textbf{Compel possibile:} se un aspetto del PG lega orgoglio o protezione di Elias alla diffidenza verso i Cacciatori, proponi che rifiuti la scorta. La complicazione è affrontare il recupero senza il loro appoggio immediato, non perdere gli indizi o provocare un attacco automatico.

\textbf{Appiglio per Nox:} i PG possono richiamarlo al compito concordato e all'ordine di contenere. Se si allarma, mostra che porta la mano all'arma e avverte; un ordine o un intervento dei PG può cambiare la situazione. Non usarlo per imporre uno scontro fuori dal loro controllo.
\end{fatebox}

\subsection{Handout}
\label{nodo-5-handout}
\phantomsection
\addcontentsline{toc}{subsubsection}{H10 -- Avviso di caccia}
\begin{handoutbox}{H10}{Avviso di caccia}
\textbf{Dove:} Affisso sulla bacheca del Rifugio o mostrato da Roland.

\textbf{Quando consegnarlo:} Quando i PG leggono la bacheca o chiedono a Roland dell'ordine; non richiede accesso a un archivio o una prova di Carisma.

\textbf{Che cosa rivela:} Elias, i sintomi riferiti, il luogo dell'intrusione e l'obiettivo di contenerlo.

\textbf{Limiti:} La competenza tecnica motiva una cautela, non certifica la trasformazione. L'avviso non è una diagnosi né un ordine di esecuzione preventiva.
\end{handoutbox}

\textbf{Senza documento:} Roland può spiegare l'ordine e riferire la testimonianza anche se i PG non prendono H10. Non consegna una mappa sostitutiva né nuovi avvistamenti non stabiliti.

\subsection{Evoluzione della scena}
\begin{clockbox}
La ricerca è già iniziata. Alla scadenza di un termine concordato senza notizie, Roland può intensificarla e chiedere conto ai PG; la scadenza non dimostra che Elias sia diventato una bestia. Una menzogna scoperta riduce la fiducia e apre un confronto, non provoca un'esecuzione.

Nox segue i PG soltanto per un incarico o una pista concreta e può perderli. Per arrivare alla sala del Progetto, i Cacciatori devono ricevere un accesso o seguire qualcuno senza perderne le tracce. Registra come ottengono il percorso. Nuove prove possono cambiare la valutazione di Roland.
\end{clockbox}

\subsection{Collegamenti}
\linkto{Ospedale Vecchio per H08 e il sorvegliante ferito; Laboratorio Alchemico per il luogo dell'intrusione; Elias e finale attraverso un percorso scoperto e un eventuale accordo con i PG. Il Rifugio non fornisce un ingresso dalle fogne.}

\nodehead{15. Nodo 6 -- Biblioteca Proibita}{Il colloquio con Elias, la runa e la lettura di un testo antico.}{Che cosa aveva compreso Elias, e che cosa suggeriscono davvero le fonti?}

\subsection{A colpo d'occhio}
\begin{overviewbox}
\begin{tabularx}{\linewidth}{>{\bfseries}p{32mm} X}
Presenti & Abel Norn riceve le richieste; Miriam Voss studia e confronta i testi.\\
Accesso & Biblioteca indipendente, legata all'Accademia dalla frequentazione degli studiosi e dalle collaborazioni di ricerca. «Proibita» è un nome tradizionale, non un divieto d'ingresso. Abel accompagna chi chiede di Miriam.\\
Handout & H11: annotazioni di Miriam sul colloquio e sul segno. H11b: preghiera pthumeriana tradotta e commentata da Maestro Willem, una delle fonti consultate.\\
Da ricordare & Miriam ricevette Elias il giorno 20 su presentazione di Silas e vide soltanto la copia della runa. Nessuno qui conosce vasca, reperto, trattamento o Celestine senza nuove informazioni portate dai PG. H11b non è una sintesi della ricerca di Willem.\\
\end{tabularx}
\end{overviewbox}

\textbf{Uso al tavolo:} segui ciò che i PG raccontano, mostrano o chiedono. I blocchi indicano intenzioni e occasioni di intervento, non una sequenza obbligatoria di domande. La lettura amplia la prospettiva sull'indagine; non è una condizione per ottenere le informazioni sulla visita di Elias.

\subsection{Arrivo e ambienti}
\subsubsection{Ingresso nella sala}
\begin{readaloud}
Fra la porta e gli scaffali rimane uno spazio libero, occupato soltanto da un tavolo dove appoggiare i volumi restituiti. Più avanti, le sedute sono raccolte intorno a tavoli larghi: un libro aperto lascia ancora posto ai fogli di chi lo consulta. Sui bordi degli scaffali sono fissati contrassegni; alcune cartelle sporgono quanto basta per leggerne il dorso.

Abel vi sente entrare, interrompe ciò che sta facendo e vi chiede chi cercate. Non vi invita a rovistare fra i libri. Quando nominate Miriam, indica il tavolo al quale lavora e vi accompagna.

La studiosa ha un volume aperto e un foglio di appunti accanto alla mano. Finisce la parola che sta scrivendo prima di alzare gli occhi. Al nome di Elias Ward posa la penna.
\end{readaloud}

\textbf{Accesso e Abel:} i PG possono chiedere di Miriam; il nome di Silas spiega il contatto, ma non è un lasciapassare obbligatorio. Essere ricevuti non permette di prelevare liberamente i volumi. Abel conserva l'ordine della raccolta e tutela la riservatezza dei lettori. Se spiegano la scomparsa, può confermare di avere udito la presentazione di Elias il giorno 20; non vide il disegno e non conosce i suoi appunti privati. La cecità non gli conferisce percezioni particolari.

\textbf{Descrizione in movimento:} usa i dettagli quando servono. Abel recupera un volume grazie ai contrassegni e alla conoscenza della raccolta; Miriam sposta i propri appunti per fare posto ai documenti dei PG. I rumori delle pagine, le sedie avvicinate e le pause per cercare un passo accompagnano azioni concrete. Non introdurre scritte mutevoli, sussurri o altri fenomeni soprannaturali.

\subsection{Interazioni e ostacoli}
\subsubsection{Parlare della visita di Elias}
\textbf{Vuole capire:} che cosa gli sia accaduto dopo il loro incontro. Se apprende della scomparsa, domanda quando sia stato visto l'ultima volta e ascolta il racconto prima di parlare della runa. Non conosce automaticamente le scoperte compiute negli altri nodi.

\textbf{Vuole comunicare:} il giorno 20 ricevette Elias su presentazione di Silas. Lui mostrò un segno copiato da un documento; lei lo confrontò con testi pthumeriani e annotò associazioni a percezione, crescita di una facoltà e possibilità di contatto. Non vide il disegno completo della vasca, il reperto o la macchina. Non conosce le azioni successive di Elias.



\textbf{Come interpretarla:} distingue ciò che ricorda dalla visita da ciò che sta apprendendo. Non attribuisce a Elias intenzioni mai dichiarate. La notizia della malattia o della ricerca dei Cacciatori le dà motivo di ascoltare, non una conoscenza improvvisa della situazione.

\textbf{Rimando agli handout:} H11 raccoglie le annotazioni del colloquio; la relativa scheda riporta le condizioni di consegna.

\subsubsection{Confrontare i documenti}
Se i PG mostrano H06b, usa questa descrizione:
\begin{readaloud}
Miriam accosta il foglio alle proprie annotazioni. Riconosce la stessa runa, poi guarda il disegno che la circonda: segue la parete della vasca e legge l'indicazione destinata al vetraio. Ruota il foglio perché possiate vederlo tutti.

«Ward mi aveva mostrato soltanto il segno. Questo disegno non l'avevo visto.»
\end{readaloud}

\textbf{Vuole chiarire:} il significato attribuito alla runa nei testi è distinto dall'uso che ne fa chi l'ha fatta incidere. H06b le rivela il supporto previsto, non il contenuto della vasca. Chiede da dove venga il documento e che cosa i PG sappiano della destinazione; non anticipa domande su esperimenti di cui nessuno le ha parlato.

\textbf{Ciò che stimola l'approfondimento:} domande su «ascolto», sul perché incidere il segno o su quanto avesse compreso Elias la portano alla fonte. Se manca H06b, può partire da H11 e dalla copia della runa: non occorre tornare all'Archivio per ascoltarla.

\begin{fatebox}
\textbf{Aspetto di scena:} \aspect{Documenti a confronto sul tavolo}, quando i fogli vengono effettivamente accostati. I PG possono confrontare dettagli, sostenere un'interpretazione o creare un vantaggio con Occulto o Erudizione. Un tiro riguarda un approfondimento specifico; non serve per leggere il disegno o ascoltare ciò che Miriam sa già spiegare. Un successo non dimostra l'efficacia della runa.
\end{fatebox}

\subsubsection{Leggere la fonte di Willem}
\textbf{Quando proporla:} Miriam dà prima una spiegazione breve, legata alla crescita della percezione. Quando i PG chiedono di precisare o mostrano interesse per ciò che il segno potrebbe significare, propone di leggere il passo su cui si basa. Non richiede una prova di fiducia e non interrompe una domanda urgente su Elias con una lezione.

\begin{readaloud}
Miriam chiede ad Abel il volume consultato con Ward. Quando lo posa sul tavolo, sposta i fogli invece di sovrapporli alle pagine. Cerca il passo, torna indietro di qualche carta e tiene aperto il volume.

Il segno non è una figura isolata. Compare accanto alle ultime righe di una preghiera. La traduzione porta il nome di Maestro Willem; sotto, due annotazioni sono separate da uno spazio. La seconda riprende e mette in dubbio una parola della prima.

Miriam vi indica l'inizio della traduzione, poi ritira la mano e vi lascia leggere.
\end{readaloud}



\textbf{Durante la lettura:} lascia spazio alle reazioni prima di spiegare. La radice che non conosce chi la calpesta, la richiesta di far crescere una facoltà e la nota sugli occhi offrono motivi diversi per fare domande. Non richiedere che i giocatori scelgano un'interpretazione prestabilita.

\textbf{Ciò che Miriam vuole comunicare, se interrogata:}
\begin{itemize}
\item La preghiera si rivolge a presenze considerate superiori, forse vicine ma non percepite. Non ne descrive con certezza natura, identità o intenzioni.
\item «Ascolto» è una scelta di traduzione discussa dallo stesso Willem, non la prova che occorrano suoni o che la runa riceva messaggi.
\item La prima nota mostra come una lettura abbia orientato la ricerca di Willem verso una facoltà dell'osservatore. La nota successiva documenta il dubbio con cui rilegge il passo dopo essersi tolto gli occhi.
\item Willem afferma di comprendere diversamente, ma ammette che potrebbe cercare una giustificazione per il proprio gesto. Il documento non dimostra che mutilarsi dia accesso a una percezione superiore.
\end{itemize}

\textbf{Come dirlo:} Miriam indica la riga su cui si fonda ogni osservazione e distingue la voce della preghiera da quella del traduttore. Non introduce una classificazione degli esseri o un riassunto della cosmologia. Il brano è una delle fonti della raccolta, non la risposta definitiva al mistero.

\subsubsection{Condividere altre scoperte}
\textbf{Aggiorna soltanto ciò che apprendono insieme:}
\begin{description}[leftmargin=5mm,style=nextline]
\item[H06b.] Miriam conosce l'incisione prevista sulla vasca; ignora che cosa contenga.
\item[H09 o testimonianze sul trattamento.] Apprende del reperto vivo, dei benefici temporanei e delle ricadute nella misura in cui vengono riferiti. Può ipotizzare che qualcuno cerchi attraverso il trattamento ciò che i testi attribuiscono alla percezione; segnala che si tratta di un'ipotesi.
\item[H16.] Può confrontare la teoria dichiarata da Celestine con i testi e i risultati. Il beneficio osservato non prova la spiegazione data dalla ricercatrice.
\end{description}

\textbf{Confidenza:} condividere prove, distinguere fatti e sospetti e rispettare i documenti le dà motivo di dedicare tempo al confronto e chiedere copie. Non è una soglia nascosta per ricevere H11. Empatia può far cogliere la sua preoccupazione per la perdita dei documenti, non rivelare una conoscenza segreta del Progetto.

\textbf{Limiti:} non identifica il proprietario del reperto, non certifica un contatto e non sa manovrare la macchina o interrompere una seduta in sicurezza. H12 rimane con Elias; Miriam non lo possiede e non ne conosce il contenuto senza che le venga mostrato.

\begin{fatebox}
\textbf{Aspetto:} \aspect{Originali da preservare}. La consultazione avviene al tavolo; Miriam e Abel aiutano a confrontare e copiare, ma non autorizzano automaticamente a portare via i volumi. Chiedere una copia o leggere H11b non richiede punti Fate.

\textbf{Compel possibile:} se un aspetto del PG esprime curiosità ostinata o ambizione di studioso, proponi che prolunghi il confronto oltre il previsto, mettendo in difficoltà un appuntamento o un termine già concordato. Specifica quale impegno rischia prima dell'accettazione; non inventare una scadenza o una punizione soprannaturale.
\end{fatebox}

\subsection{Handout}
\label{nodo-6-handout}
\phantomsection
\addcontentsline{toc}{subsubsection}{H11 -- Note sulla runa}
\begin{handoutbox}{H11}{Note sulla runa}
\textbf{Dove:} Le annotazioni di Miriam sul colloquio del giorno 20.

\textbf{Quando consegnarlo:} Quando i PG chiedono del colloquio, mostrano il segno o vogliono sapere che cosa Miriam avesse spiegato a Elias. Non richiede una prova di buone intenzioni.

\textbf{Che cosa rivela:} Il resoconto della lettura della runa al giorno 20: associazioni a percezione, crescita di una facoltà e possibilità di contatto.

\textbf{Limiti:} Miriam vide soltanto la copia del segno. Non aggiornare H11 con informazioni apprese oggi; la lettura non dimostra un contatto.
\end{handoutbox}

\phantomsection
\addcontentsline{toc}{subsubsection}{H11b -- Preghiera tradotta e commentata da Willem}
\begin{handoutbox}{H11b}{Preghiera tradotta e commentata da Willem}
\textbf{Dove:} Il passo si trova nel volume consultabile della Biblioteca; Abel sa recuperarlo.

\textbf{Quando consegnarlo:} Durante la lettura proposta da Miriam, quando i PG chiedono di precisare il significato del segno o mostrano interesse. Se Miriam manca, Abel può far consultare il volume già tradotto, senza sostituirla nell'interpretazione.

\textbf{Che cosa rivela:} Una preghiera tradotta e commentata da Maestro Willem. La runa accanto all'ultima invocazione è la stessa di H11 e H06b, distinta dal marchio amministrativo. Il foglio riproduce il passo e i commenti presenti nel volume.

\textbf{Limiti:} È una delle fonti della raccolta, non un documento proveniente dal Progetto o contenente appunti di Elias. Gli originali restano soggetti alle condizioni di consultazione della Biblioteca.
\end{handoutbox}

\begin{gmbox}
H11b è un testo originale della campagna attribuito a Willem nella finzione. Non è una citazione del gioco. Non introduce una procedura rituale, una cura, un'entità identificata o un effetto soprannaturale verificato.
\end{gmbox}

\subsection{Evoluzione della scena}
Abel può confermare la visita e far consultare il volume con H11b, già tradotto; non occorre tradurre dal pthumeriano per leggerlo. Sa ritrovarlo, non sostituire Miriam nell'interpretazione. Per confronti ulteriori serve una competenza pertinente o un nuovo colloquio con lei.

Se i PG vogliono soltanto ricostruire la visita, Miriam racconta il colloquio e mostra H11 senza imporre la lettura di H11b. Se non prendono gli handout, possono comunque leggere e discutere i testi. Non perdono un accesso fisico al finale: la Biblioteca non custodisce una pianta delle catacombe o un nuovo percorso verso il Progetto.

\begin{clockbox}
Con la sola runa, Miriam non sa quale luogo cercare. Se i PG le rivelano reperto e trattamento, può chiedere di studiare le prove o proporre di accompagnarli. Per raggiungere la sala deve ricevere un percorso utilizzabile e ottenere accesso, oppure essere accompagnata. Non arriva al finale per il solo trascorrere del tempo.

Un eventuale interesse per campioni riguarda materiale già disponibile e informazioni effettivamente ricevute; non è un incarico obbligatorio o un pedaggio per spiegare la runa.
\end{clockbox}

\subsection{Collegamenti}
\linkto{Archivio Ecclesiastico per il documento da cui Elias copiò la runa; Ospedale e Laboratorio Alchemico soltanto attraverso destinazioni e informazioni emerse dai documenti o dai racconti dei PG. La Biblioteca offre interpretazioni e collaborazione, non un accesso sotterraneo.}

\nodehead{16. Nodo 7 -- Laboratorio Alchemico}{Verifica delle consegne e scelta di un accesso all'ala inferiore.}{Dove è stato trasferito il regolatore e come possono seguirne il percorso?}

\subsection{A colpo d'occhio}
\begin{overviewbox}
\begin{tabularx}{\linewidth}{>{\bfseries}p{32mm} X}
Presenti & Addetto all'ascensore e sorvegliante. Il tecnico della macchina è presente se impegnato su forniture o manutenzione; il personale ordinario non conosce automaticamente gli esperimenti.\\
Accesso & Ingresso principale sulla sala di lavoro; porta esterna sul fianco direttamente nel deposito, aperta dalla chiave di Silas. La porta sorvegliata sul fondo e l'ascensore restano controlli distinti.\\
Handout & H13 aperto sul piano di registrazione vicino alla porta di servizio; H14 affisso nel settore controllato dell'ascensore, al piano del deposito.\\
Da ricordare & Il Laboratorio svolge preparazioni alchemiche per l'Ospedale; il deposito riceve anche le forniture del Progetto. H13 collega la commessa 44-14 all'ala inferiore. H12 rimane con Elias: non si trova qui e non serve a raggiungerlo.\\
Vie da valutare & Scala dell'Ospedale, ascensore e galleria, scarico dalle fogne. H14 ne mostra i collegamenti, non tutta la rete fognaria cittadina.\\
\end{tabularx}
\end{overviewbox}

\textbf{Uso al tavolo:} scegli il blocco corrispondente alle azioni dei PG. Una normale consegna ospedaliera non prova un abuso: il nodo serve a rintracciare il regolatore e verificare il tratto successivo del suo viaggio. Nessuna prova unica risolve ingresso, documenti e discesa.

\subsection{Arrivo e ambienti}
\subsubsection{Disposizione e ritmo del luogo}
L'ingresso principale sulla strada apre sulla sala di lavoro. Il passaggio sul fondo conduce al deposito: entrando dalla sala, gli scaffali sono a sinistra e la porta esterna di servizio a destra. Questa porta dà sul passaggio lungo il fianco dell'edificio ed è quella aperta dalla chiave di Silas. Il piano di registrazione è vicino alla porta di servizio. Sul fondo del deposito, oltre la corsia dei carrelli, la porta sorvegliata conduce al pianerottolo dell'ascensore.

\textbf{Per descrivere la scena:} adatta attività e presenze all'ora della visita. Durante un carico, mostra gli addetti che controllano, registrano e spostano le casse. Durante una pausa, indica dove si trovano persone e materiali. Una consegna parte quando è prevista nella scena, non automaticamente all'arrivo dei PG. Recipienti, imballaggi e rumori sono dettagli del luogo: diventano rilevanti se i giocatori decidono di usarli. Richiedi una prova soltanto per un'azione incerta con un ostacolo concreto.

\subsubsection{Dalla strada}
\begin{readaloud}
Il Laboratorio occupa l'edificio accanto all'Ospedale. L'ingresso principale si apre sulla strada; attraverso le finestre si distinguono persone chine sui tavoli di lavoro. Sul fianco dell'edificio corre un passaggio di servizio. Una porta larga, con la soglia consumata e il legno segnato all'altezza delle casse, dà accesso al deposito.

Dall'interno arrivano il tintinnio dei recipienti e il rumore intermittente di qualcosa trascinato sul pavimento.
\end{readaloud}

\textbf{Orientamento:} i PG possono avvicinarsi all'ingresso principale, osservare dalle finestre o percorrere il fianco fino alla porta di servizio. La serratura esterna è quella per cui Silas possiede la chiave. Indica che cosa è effettivamente visibile dalla posizione scelta; il deposito può essere occupato anche quando la porta esterna è chiusa.

\subsubsection{Sala di lavoro}
\begin{readaloud}
Due tavoli lunghi occupano il centro della sala, lasciando un corridoio fra l'ingresso e il passaggio sul fondo. Sui piani sono disposti mortai, bilance, recipienti di vetro e piccoli vassoi di strumenti. Gli scaffali lungo le pareti raccolgono bottiglie e contenitori chiusi; accanto a una vasca di lavaggio, panni umidi sono stesi sul bordo.

Un lavorante versa una preparazione attraverso un imbuto mentre un altro pesa il contenuto di un recipiente. Chi termina un'operazione libera una porzione del tavolo e vi dispone gli strumenti per la successiva. Dal passaggio sul fondo si vedono scaffalature più basse e il fianco di una cassa.

Quando entrate, il lavorante più vicino alza gli occhi e vi chiede di chi avete bisogno.
\end{readaloud}

\subsubsection{Deposito}
\begin{readaloud}
Entrando dalla sala di lavoro, avete gli scaffali sulla sinistra e la porta esterna di servizio sulla destra. Le casse più ingombranti sono appoggiate a terra lungo le pareti. Fra loro rimane una corsia abbastanza larga da far passare un carrello carico. Spago tagliato, listelli e materiale d'imballaggio sono raccolti vicino alla porta esterna.

Un piano di registrazione sostiene un inventario aperto, una penna e alcune ricevute tenute ferme da un peso. Le colonne dell'inventario distinguono materiale ricevuto, quantità presenti e trasferimenti. Accanto, un carrello attende con il carico sistemato sulle assi. Sul fondo della corsia si trova la porta del settore merci. Il sorvegliante occupa una sedia accanto allo stipite, dalla quale può vedere il passaggio interno e la porta di servizio.
\end{readaloud}

\textbf{Da un altro ingresso:} se i PG entrano dalla porta di servizio, descrivi il piano vicino a loro, le scaffalature di fronte e i due passaggi interni: quello verso la sala di lavoro e quello controllato verso l'ascensore. La visuale della guardia dipende dalla posizione effettiva e dagli eventuali carichi in movimento; non trattare le casse come un nascondiglio garantito.

\subsubsection{Pianerottolo dell'ascensore}
\begin{readaloud}
Oltre la porta, il pavimento continua fino alla soglia dell'ascensore. Davanti alla cabina rimane uno spazio di manovra segnato dalle ruote dei carrelli. Il comando è fissato alla parete, vicino alla posizione dell'addetto; poco più in là è appeso lo schema di servizio. Nel disegno si distinguono la scala dell'Ospedale, la galleria delle forniture e lo scarico verso le fogne. Sotto sono riportate le avvertenze per il lavaggio.
\end{readaloud}

\textbf{Quando l'ascensore si muove:} il rumore del meccanismo riempie il pianerottolo e si sente anche dal deposito. L'addetto segue il movimento della cabina, poi torna al carico in attesa. Usa questa descrizione durante un trasferimento effettivo: il rumore segue il funzionamento del mezzo.

\subsection{Interazioni e ostacoli}
\subsubsection{Chiedere di una consegna}
\textbf{Richieste ordinarie:} i PG possono chiedere di verificare una consegna, citando H05, H06a o la commessa 44-14. Il personale può indirizzarli a chi gestisce i carichi; una richiesta non autorizza a circolare nel deposito o scendere all'ala inferiore. Non serve fingere che tutto l'edificio sia segreto.

\textbf{Aspetto:} \aspect{Laboratorio in attività}. Il lavoro offre occasioni per parlare con gli addetti o osservare i movimenti; una distrazione deve coinvolgere persone e operazioni presenti nella scena. Non introdurre incidenti alchemici o conoscenze degli esperimenti solo per generare un indizio.

\subsubsection{Verificare la consegna nel deposito}
\textbf{Documento visibile:} nomina l'inventario aperto sul piano; il codice H13 serve soltanto al GM. Da lontano si riconosce la funzione delle colonne, mentre per leggere le singole voci occorre avvicinarsi.

\textbf{Mostrare l'uso durante una consegna:} l'addetto legge il riferimento sul collo, torna al piano e scorre l'inventario con un dito. Annota qualcosa, poi indica dove sistemare la cassa. Questo gesto rende visibile il legame fra merci e registrazione. Se il lavoro è fermo, basta descrivere il documento e le sue colonne: non introdurre un carico soltanto per attirare l'attenzione.

\textbf{Consultazione:} se i PG chiedono della commessa 44-14 o del regolatore, l'addetto consulta H13 e può mostrare la voce pertinente come normale verifica della consegna. La richiesta non richiede automaticamente Carisma. Se preferiscono leggere di nascosto, chiarisci come raggiungono il piano, chi li vede e quanto possono soffermarsi; un'eventuale prova riguarda quell'azione.

\subsubsection{Superare la porta controllata}
\textbf{Addetto all'ascensore:} conosce carichi, inventario, galleria e destinazione delle forniture. Custodisce la chiave che abilita l'ascensore. Può indicare le operazioni previste in quel momento; non esiste un orario universale di trasferimento ricavabile da H13.

\textbf{Sorvegliante:} è diverso dall'uomo ferito da Elias. Controlla permessi e passaggi riservati, sa dell'intrusione e deve riferire altri accessi non autorizzati. Una richiesta di verifica delle consegne non autorizza a scendere. Un ordine di Aldren vale per ciò che autorizza espressamente; per l'accesso alla sala del Progetto occorre la conferma già prevista dalla sua scheda.

\textbf{Tecnico della macchina:} se presente per un carico o un intervento, sa che il regolatore stabilizza la pressione. Non sa produrre un contatto e non condivide queste conoscenze con tutto il personale del Laboratorio. Motiva la presenza con il suo lavoro, senza aggiungere un nuovo PNG.

\begin{fatebox}
\textbf{Aspetto:} \aspect{Accesso merci controllato}. I PG possono chiedere un permesso, sostenere una richiesta concreta con Carisma, presentare un pretesto con Inganno o preparare una distrazione. Furtività riguarda il passaggio davanti a una persona o il raggiungimento di un oggetto specifico, non l'intero ingresso nel Progetto.

\textbf{Chiave di Silas:} apre la porta esterna sul fianco dell'edificio e permette di entrare direttamente nel deposito. La porta controllata e l'ascensore restano passaggi distinti, con il personale presente.

\textbf{Compel possibile:} un aspetto pertinente alla protezione di Elias o all'insofferenza per l'autorità può portare un PG a contestare apertamente la guardia, rendendo riconoscibile l'interesse del gruppo per l'intruso. Concorda la complicazione; non far arrivare rinforzi senza una comunicazione concreta.
\end{fatebox}

\subsubsection{Usare l'ascensore}
\textbf{Uso dell'ascensore:} la chiave ne abilita il funzionamento; non deve essere tenuta da chi viaggia. In condizioni ordinarie l'addetto lo abilita per il carico autorizzato. I PG devono ottenere la collaborazione necessaria o descrivere come superano quel controllo: una competenza tecnica può aiutare in un tentativo concreto, non fornisce automaticamente la chiave.

\textbf{Scelta dei PG:} consultare H14 non li obbliga a usare l'ascensore. Possono valutare la scala dell'Ospedale o cercare lo scarico nelle fogne.

\textbf{Rimando agli handout:} la scheda H14 raccoglie le condizioni di consultazione e i collegamenti mostrati dallo schema.

\subsubsection{Ricostruire l'ingresso di Elias}
L'addetto aveva già aperto l'accesso e abilitato l'ascensore per un trasferimento. Elias sfruttò il movimento delle forniture per raggiungere il settore merci. Il sorvegliante cercò di trattenerlo; Elias lo ferì al braccio per liberarsi e lo risparmiò. L'addetto si allontanò per chiamare aiuto. Rimasto momentaneamente senza osservatori, Elias raggiunse la cabina e scese con l'ascensore già abilitato, senza rubare la chiave. Né il sorvegliante né l'addetto videro la discesa. L'allarme interno partì subito e raggiunse Celestine; segnalava un intruso armato di cui si erano perse le tracce nel deposito. Durante i soccorsi e i primi controlli Elias percorse la galleria, attraversò la sala del Progetto riconoscendo il regolatore e si nascose nell'intercapedine del vano di manutenzione. Il giorno 27 Celestine fece coinvolgere i Cacciatori.

\textbf{Fonti per i PG:} il sorvegliante ferito, rintracciabile all'Ospedale attraverso Roland, racconta lo scontro e ammette di aver perso di vista Elias. L'addetto riferisce di essersi allontanato per chiamare aiuto lasciando l'ascensore abilitato. Entrambi possono sospettare che Elias sia sceso, ma nessuno dei due lo ha visto farlo o sa dove si trovi. Il percorso descritto sopra è un fatto noto al GM. Elias ha conosciuto gli altri collegamenti dallo schema: non ha percorso le fogne per entrare.

\textbf{Dopo l'ingresso:} Elias trova riparo nell'intercapedine del vano di manutenzione. Rimane per preparare un sabotaggio del regolatore e comprendere il reperto; posizione, limiti del riparo ed effetti dell'intervento sono descritti nel nodo finale.

\subsubsection{Raggiungere lo scarico dalle fogne}
\textbf{Raggiungere l'imbocco:} le fogne attraversano la città. H02 resta una pianta generale; H14 documenta il collegamento con lo scarico ospedaliero, non un itinerario dalla strada. Il GM sceglie gli ingressi cittadini raggiungibili, il percorso e le difficoltà coerenti con la scena. Distingui trovare l'imbocco dello scarico da superarne la grata: le prove non devono mettere in dubbio l'esistenza del collegamento già documentato.

\textbf{Direzione:} l'acqua defluisce dal vano di manutenzione dell'Ospedale verso le fogne. Per entrare i PG percorrono lo scarico in direzione opposta. Fuori dal lavaggio è abbastanza ampio da attraversare con difficoltà; non è un passaggio ordinario del personale.

\textbf{Grata:} normalmente chiusa e fissata in posizione. L'acqua passa fra le sbarre; gli addetti operano dal lato interno e non la aprono a ogni lavaggio. Non c'è una nuova chiave da cercare. Rimuoverla o forzarla richiede un tentativo concreto: valuta attrezzi, tempo e rumore. Non trattarla come invalicabile, ma non basta trovare lo scarico per oltrepassarla.

\textbf{Lavaggio:} comincia alla quarta campana notturna e dura circa un'ora. Il flusso rende il passaggio impraticabile. H14 riporta orario e direzione dell'acqua; quando i PG pianificano, chiarisci quanto manca rispetto all'ora della scena. Non è un codice da decifrare e non richiede segnali a colpi. H13 non stabilisce l'orario del lavaggio o di ogni carico.

\textbf{Arrivo:} oltre la grata si raggiunge il vano di manutenzione adiacente alla sala del Progetto, nell'ala inferiore. Non si entra nel Laboratorio Alchemico e non si aprono i suoi accessi controllati.

\subsection{Handout}
\label{nodo-7-handout}
\phantomsection
\addcontentsline{toc}{subsubsection}{H13 -- Inventario del deposito}
\begin{handoutbox}{H13}{Inventario del deposito}
\textbf{Dove:} Aperto sul piano di registrazione vicino alla porta esterna di servizio.

\textbf{Quando consegnarlo:} Quando i PG riescono a consultarlo tramite l'addetto o raggiungendo le carte. L'esistenza del documento è visibile: non occorre una prova di Indagine per scoprirla. L'ingresso nel deposito, da solo, non equivale ad averne letto il contenuto.

\textbf{Che cosa rivela:} Il regolatore della commessa 44-14 è stato trasferito all'ala inferiore riservata. Il numero collega l'inventario ai documenti della Fonderia e dell'Archivio. Le altre voci mostrano forniture in transito e consumo ordinario. Il rimando a H14 segnala uno schema dei collegamenti consultabile presso l'ascensore.

\textbf{Limiti:} Non indica il peso del solo regolatore. Inviare materiali all'Ospedale è normale: il documento non dimostra da solo coercizione, finalità ultraterrene o complicità degli addetti. Tubi e strumenti generici non permettono di dedurre con Tecnica il funzionamento della macchina nascosta.
\end{handoutbox}

\phantomsection
\addcontentsline{toc}{subsubsection}{H14 -- Schema degli accessi}
\begin{handoutbox}{H14}{Schema degli accessi}
\textbf{Dove:} Affisso nel settore controllato dell'ascensore, al piano del deposito. H13 ne indica l'esistenza.

\textbf{Quando consegnarlo:} Quando i PG raggiungono il settore e consultano lo schema, prima di scendere. Un addetto può mostrarlo durante una consultazione autorizzata. Non richiede di trovare un pannello nascosto o un'altra sala macchine.

\textbf{Che cosa rivela:} Ascensore e galleria verso il corridoio dell'ala inferiore, scala ospedaliera verso le stanze dei pazienti, scarico verso il vano di manutenzione. Riporta grata, direzione del deflusso e lavaggio alla quarta campana notturna. «Sala trattamenti» indica la sala del Progetto, sul medesimo livello delle stanze.

\textbf{Limiti:} Non contiene annotazioni di Elias, non spiega il trattamento e non descrive tutta la rete fognaria cittadina. Consultarlo non obbliga a usare l'ascensore.
\end{handoutbox}

\textbf{Senza H13 o H14:} gli addetti possono confermare il trasferimento del regolatore e spiegare o mostrare i collegamenti che conoscono, se le circostanze del colloquio lo permettono. Non consegnano automaticamente ogni informazione a chi entra. Le prove scritte permettono confronti e pianificazione, ma non sono una sequenza obbligatoria.

\textbf{H12 -- Scoperte e dubbi di Elias:} rimane in suo possesso e può essere mostrato al successivo incontro. Non si trova nel deposito, non contiene istruzioni rivolte ai PG e non è necessario per raggiungere il finale.

\subsection{Evoluzione della scena}
\begin{clockbox}
I controlli tengono conto dell'intrusione già avvenuta. Ulteriori restrizioni dipendono da eventi osservati o riferiti, non dalla fine della scena. Se il sorvegliante chiama aiuto, stabilisci come il messaggio raggiunge il destinatario e lascia ai PG modo di reagire. Altre fazioni arrivano soltanto per un motivo concreto. Solo una minaccia giunta a Celestine può innescare le misure sui materiali previste dall'orologio.
\end{clockbox}

\subsection{Collegamenti}
\linkto{Ala inferiore tramite ascensore e galleria, scala ospedaliera o scarico dalle fogne. Fonderia e Archivio per il confronto della commessa 44-14; Ospedale e Cacciatori per la testimonianza del sorvegliante ferito.}

\nodehead{17. Nodo finale -- Il Progetto}{Incontro con Elias, interruzione degli abusi e destino delle prove.}{Che cosa scelgono di proteggere i PG, e come intervengono nella situazione in corso?}

\subsection{A colpo d'occhio}
\begin{overviewbox}
\begin{tabularx}{\linewidth}{>{\bfseries}p{32mm} X}
Presenti & Celestine e il tecnico nella sala trattamenti; sei pazienti e personale nelle stanze; Elias nell'intercapedine del vano di manutenzione. Alleati soltanto se accompagnano i PG o hanno ottenuto concretamente un accesso.\\
Accesso & Scala ospedaliera, galleria delle forniture oppure scarico verso il vano di manutenzione, secondo il percorso ottenuto dai PG.\\
Handout & H12 con Elias; H15 alla postazione del tecnico; H16 fra le carte di lavoro di Celestine.\\
Da ricordare & Si comincia fra due sedute: la macchina mantiene le condizioni del reperto e nessun paziente è collegato. Fra i sei nelle stanze ci sono i due cancellati dai registri.\\
Attività in corso & Celestine confronta risultati e prepara il lavoro successivo; il tecnico controlla macchina e collegamenti; il personale assiste e accompagna i pazienti. Elias prepara un intervento sul regolatore e cerca ancora di capire il reperto.\\
\end{tabularx}
\end{overviewbox}

\textbf{Preparazione:} questa è la situazione di base. Aggiornala se i PG hanno provocato un allarme, ottenuto un'ispezione o indotto Celestine a preparare un trasferimento. Segna posizione degli addetti, controlli già attivi e percorso degli eventuali alleati. L'orologio modifica le condizioni, non impone un incontro o un esito.

\textbf{Ritmo:} la seduta successiva richiede di andare a prendere un paziente, accompagnarlo e prepararlo. Mostra queste azioni se avvengono, lasciando ai PG spazio per intervenire. Una pausa nella conversazione al tavolo non fa partire automaticamente il trattamento.

\subsection{Arrivo e ambienti}
\subsubsection{Disposizione dell'ala e arrivo}
Le stanze dei pazienti si aprono sul corridoio raggiunto dalla scala ospedaliera. La galleria dell'ascensore confluisce nello stesso corridoio; da lì si accede alla sala trattamenti, indicata così in H14 e chiamata sala del Progetto nel manuale. Il vano di manutenzione è adiacente alla sala e comunica con essa; lo scarico delle fogne termina in quel vano. Tutti questi ambienti appartengono al medesimo livello.

\begin{description}[leftmargin=5mm,style=nextline]
\item[Dalla scala.] I PG raggiungono prima le stanze dei pazienti e il personale che le serve. Applica il permesso o le conseguenze con cui hanno superato il controllo ospedaliero.
\item[Dalla galleria.] Arrivano nel corridoio delle forniture e vedono gli accessi alle stanze e alla sala. Un accompagnamento dell'addetto resta limitato all'incarico concordato.
\item[Dallo scarico.] Superati grata e percorso, entrano nel vano di manutenzione vicino al riparo di Elias. Lo scarico non dà accesso al Laboratorio Alchemico.
\end{description}

\subsubsection{Corridoio e stanze dei pazienti}
\begin{readaloud}
Le porte delle stanze si aprono lungo il corridoio. Attraverso quelle aperte vedete letti, coperte ripiegate e recipienti d'acqua lasciati a portata di mano. Il personale si ferma accanto ai letti per parlare con i degenti e aiutarli a cambiare posizione. Alcune risposte arrivano subito; altre richiedono di ripetere la domanda.

Più avanti, il rumore regolare delle pompe proviene dalla sala trattamenti. Chi si affaccia dalle stanze può vedere le persone che vengono accompagnate in quella direzione.
\end{readaloud}

\textbf{I sei:} mantieni il numero e le identità già riportate da Maud. La cancellazione di due nomi riguarda i registri, non una morte o un'ulteriore sparizione. Prima della scena annota chi riesce a camminare e chi ha bisogno di un braccio o di assistenza continuativa: almeno alcuni si muovono autonomamente, altri sono troppo deboli per farlo da soli. Mantieni queste condizioni nel seguito, senza introdurre nuove diagnosi.

\textbf{Presenza concreta:} un degente può chiedere se siano venuti per riportarlo di sopra; un altro può rifiutare di alzarsi quando riconosce i preparativi di una seduta. Se Oren o Maud sono presenti, i pazienti possono riconoscerli. Queste reazioni mostrano le conseguenze già documentate, non aggiungono una nuova pista.

\textbf{Muoversi e aiutare:} il personale può spiegare quali degenti abbiano bisogno di sostegno. Per portarli via occorre organizzare accompagnatori e una via accessibile. Oren e Maud, se presenti, offrono assistenza; non sono componenti obbligatorie di una soluzione tecnica. La sospensione delle sedute imposte non rende permanenti i miglioramenti né cancella le loro malattie.

\subsubsection{Sala trattamenti}
\begin{readaloud}
Il lettino occupa lo spazio di lavoro accanto alla macchina. Le condotte seguono il bordo dell'impianto fino alla vasca; sul vetro riconoscete la runa del disegno destinato al vetraio. Dentro la vasca il reperto organico rimane immerso nel sistema di conservazione.

Il tecnico controlla i collegamenti dalla propria postazione. Uno schema è disposto fra gli strumenti di lavoro. Sul tavolo di Celestine sono aperte le registrazioni delle sedute, con fogli separati per osservazioni e disposizioni. Lei confronta due annotazioni, poi torna a rivolgersi al tecnico.

Fra il lettino, la postazione e il tavolo restano i passaggi usati per lavorare e accompagnare il paziente. Una porta laterale conduce al vano di manutenzione, dove proseguono le condotte.
\end{readaloud}

\textbf{Orientamento:} colloca ingresso dal corridoio, porta del vano, lettino e postazioni quando imposti la scena e mantieni quelle posizioni. Mostra chi può vedere un avvicinamento al regolatore o alle carte. Strumenti, fogli e passaggi sono elementi del luogo; diventano appigli se i PG li usano, senza una prova obbligatoria per ogni dettaglio.

\textbf{Macchina:} vasca, pompe e condotte formano l'impianto. Durante una seduta il sangue esce dal paziente, passa attraverso il sistema con il reperto e torna al paziente. Il regolatore aggiunto stabilizza la pressione e permette una portata maggiore. Il confronto con H03 o la testimonianza di Elias ne confermano l'origine. La runa e i benefici temporanei non dimostrano un contatto.

\subsubsection{Vano di manutenzione e riparo}
\begin{readaloud}
Nel vano, le condotte seguono la parete della sala trattamenti e lasciano uno spazio stretto per gli interventi. Lo scarico prosegue verso le fogne; dalla grata arriva il rumore dell'acqua nel collettore. Lungo la parete tecnica, pannelli rimovibili danno accesso al retro delle tubazioni.

Le voci della sala arrivano attraverso la porta quando è aperta; il funzionamento dell'impianto copre parte dei suoni più bassi.
\end{readaloud}

\textbf{Intercapedine:} dietro un pannello c'è uno spazio basso lungo pochi metri, sufficiente a rannicchiarsi. Elias lo ha aperto con i propri attrezzi dopo l'ingresso e ha riaccostato il pannello dall'interno, conservando le viti rimosse. Non è una quarta entrata e non porta ad altre stanze. Chi esamina quel pannello può notarne il fissaggio mancante; la ricerca approfondita e il rumore di un movimento possono esporre il riparo.

\textbf{Permanenza:} acqua e provviste portate da casa coprono l'attesa iniziale. Le prime ispezioni si sono concentrate sui locali praticabili del complesso: nessuno ha visto Elias scendere e le ricerche sono distribuite fra il deposito e i possibili passaggi. Il personale conosce lo spazio delle condotte ma non lo ha ancora aperto nella ricerca dell'intruso. Elias rischia di essere visto quando esce. Se sono trascorsi altri giorni, considera risorse consumate e ispezioni effettive: il riparo non lo protegge indefinitamente.

\subsection{Interazioni e ostacoli}
\subsubsection{Parlare con Elias}
Vuole impedire altri abusi senza morire nel tentativo. Prepara il sabotaggio del regolatore fra due sedute, ma cerca ancora di comprendere il reperto e la possibilità di un contatto. L'interesse è personale e intellettuale: non sente una chiamata soprannaturale. La curiosità contribuisce a rimandare un intervento definitivo.

\textbf{Lucidità:} perde il filo, ricontrolla operazioni note e a volte riconosce di aver reagito con troppa irritazione. Può ragionare, scegliere e rivedere il piano. La malattia precede l'ingresso; fermare il Progetto non lo guarisce.

\textbf{Occasioni d'incontro:} dal vano può sentire chi entra lì. Dal riparo coglie soltanto ciò che arriva attraverso porta, distanza e rumore: non ascolta tutte le conversazioni dell'ala. Se i PG si avvicinano alla sala o parlano con Celestine, può accorgersi di loro e scegliere di intervenire. Sentire il nome di Eleanor o una proposta di fermare le sedute gli offre un motivo per esporsi; non lo fa comparire automaticamente.

\textbf{Informazioni e alternative:} riconosce il regolatore e ne spiega la funzione. Ammette di aver risparmiato il sorvegliante. I PG possono offrirgli una diversa possibilità di fermare gli abusi, preservare le prove o uscire. Non sa in anticipo che siano stati chiamati da Eleanor.

\subsubsection{Intervenire sul regolatore}
\textbf{Obiettivo immediato:} raggiungere il regolatore nell'intervallo fra le sedute e renderlo inutilizzabile. Richiede un'occasione e un intervento concreto sul componente; non è già avvenuto al momento iniziale e non si completa fuori scena soltanto perché i PG arrivano.

\textbf{Se riesce:} il tecnico sospende la preparazione del trattamento alla portata aumentata e deve intervenire sul guasto. Il reperto resta nella vasca e i sei pazienti, inizialmente nelle stanze, non subiscono un danno automatico. Il sabotaggio del solo regolatore non equivale alla distruzione delle funzioni di conservazione dell'impianto.

\textbf{Limite:} le prime prove avvenivano senza quel componente e la cerchia possiede già i disegni. Il regolatore può essere ricostruito; H12 contiene osservazioni, non il progetto costruttivo necessario a farlo. Fermare gli abusi stabilmente richiede anche una decisione sulla custodia dei pazienti e sull'autorità che dispone le sedute.

\textbf{Azioni distinte:} impedire una nuova seduta, intervenire su una seduta già avviata e distruggere l'impianto sono situazioni diverse. Se il gioco porta a un paziente collegato, considera quello stato concreto; non introdurre una procedura segreta di scollegamento sicuro o una morte automatica allo spegnimento. La situazione iniziale fra le sedute evita di imporre questo problema all'arrivo.

\subsubsection{Affrontare Celestine}
\textbf{Vuole difendere:} i miglioramenti reali osservati e la possibilità di renderli duraturi. Confronta i referti per preparare la seduta successiva. Quando difende la ricerca, mostra un'osservazione favorevole e ricorda un paziente che ha riconosciuto un familiare. La durata del beneficio e le ricadute emergono confrontando carte, testimonianze e H09.

\textbf{Arrivo dei PG:} davanti a una visita autorizzata chiede l'oggetto della verifica e cerca di controllarla. Davanti a persone scoperte nell'area riservata pretende un'identificazione e chiama il personale. Una minaccia al reperto o alla macchina la spinge a impedirla e chiedere aiuto. Ogni allarme deve essere comunicato concretamente.

\textbf{Davanti alle prove:} può ammettere le ricadute e sostenere proprio per questo più durata, frequenza o portata. Sa che i rifiuti vengono superati e i rapporti esterni omettono gli esiti sfavorevoli. La sua interpretazione del contatto resta un'ipotesi, anche quando la difende con convinzione.

\textbf{Offerta già disponibile:} riservatezza in cambio di accesso controllato ai dati e assistenza a Elias. Se propone di sottoporlo al trattamento, i benefici osservati rimangono temporanei: non possiede una cura definitiva.

\textbf{Fermarla:} una sospensione può derivare da confronto, perdita del sostegno del personale o di Aldren, allontanamento dei pazienti o impossibilità materiale di usare l'impianto. Non richiede che Celestine riconosca immediatamente il proprio errore. Il tecnico conosce il funzionamento, il personale vede i pazienti: ciascuno reagisce a ciò che sa, senza condividere automaticamente la teoria della direttrice.

\subsubsection{Coinvolgere gli alleati presenti}
\begin{tabularx}{\linewidth}{>{\bfseries}p{29mm} X}
Roland e Nox & Proteggono uscite e presenti secondo gli accordi. Aspetto alterato e spasmi non autorizzano un'esecuzione; intervengono davanti a un'aggressione incontrollata non contenibile.\\
Oren e Maud & Riconoscono e assistono i pazienti; aiutano a organizzare gli spostamenti e testimoniano ciò che hanno osservato.\\
Aldren & Può mettere in discussione autorizzazioni e sostegno alla ricerca quando apprende gli abusi. Chiede custodia delle prove; non comanda automaticamente tecnici e Cacciatori.\\
Miriam & Aiuta a conservare e interpretare documenti. Può proporre contatti con studiosi dell'Accademia, non impegnare l'istituzione o trasferirle il Progetto per propria autorità.\\
\end{tabularx}

\textbf{Accesso e conoscenza:} includi soltanto chi è arrivato attraverso un percorso ottenuto o accompagnando i PG. Aggiorna ciò che ciascuno sa in base a quanto ha visto o ricevuto. Nessuna fazione compare per la sola apertura della scena finale.

\begin{fatebox}
Scegli sfida, contesa o conflitto in base alle azioni effettive. Eventuali aspetti di scena devono corrispondere a fatti presenti: \aspect{Ricerca senza consenso}, \aspect{Documenti sul tavolo di Celestine}. Pericoli, vantaggi e opposizioni dipendono da posizioni e persone descritte; non richiedere una prova per ogni dettaglio d'ambiente. Questa situazione non stabilisce una sequenza obbligatoria di interventi.
\end{fatebox}

\subsection{Handout}
\label{nodo-8-handout}
\phantomsection
\addcontentsline{toc}{subsubsection}{H12 -- Scoperte e dubbi di Elias}
\begin{handoutbox}{H12}{Scoperte e dubbi di Elias}
\textbf{Dove:} Con Elias.

\textbf{Quando consegnarlo:} Può mostrarlo durante il colloquio.

\textbf{Che cosa rivela:} Osservazioni e conflitto fra fermare la ricerca e comprenderla.

\textbf{Limiti:} Rimane facoltativo: non è necessario per raggiungere Elias e non è il progetto costruttivo del regolatore.
\end{handoutbox}

\phantomsection
\addcontentsline{toc}{subsubsection}{H15 -- Schema della macchina}
\begin{handoutbox}{H15}{Schema della macchina}
\textbf{Dove:} Presso la postazione del tecnico.

\textbf{Quando consegnarlo:} Quando i PG riescono a consultarlo.

\textbf{Che cosa rivela:} Il circuito del trattamento e la funzione del regolatore, in un documento di servizio.

\textbf{Limiti:} Non contiene annotazioni di Elias e non prova un contatto.
\end{handoutbox}

\phantomsection
\addcontentsline{toc}{subsubsection}{H16 -- Rapporto sul Progetto}
\begin{handoutbox}{H16}{Rapporto sul Progetto}
\textbf{Dove:} Fra le carte di lavoro di Celestine.

\textbf{Quando consegnarlo:} Può essere mostrato durante una verifica o consultato raggiungendo quelle carte.

\textbf{Che cosa rivela:} Osservazioni, interpretazione e disposizioni sul Progetto.

\textbf{Limiti:} Interrogare Celestine non mette automaticamente il documento nelle mani dei PG; la sua interpretazione rimane distinta dai fatti osservati.
\end{handoutbox}

Le testimonianze, la macchina e le condizioni dei pazienti restano fonti anche senza i fogli. Non imporre una lettura prima delle decisioni dei giocatori.

\subsection{Evoluzione della scena}
\subsubsection{Conseguenze da registrare}
\begin{tabularx}{\linewidth}{>{\bfseries}p{29mm} X}
Elias & Esce, rimane, viene affidato a qualcuno, catturato o muore secondo gli eventi. Sopravvivere non risolve la sua malattia.\\
Pazienti & Chi viene sottratto alle sedute imposte, chi resta sotto controllo e chi si occupa dell'assistenza. Conta tutti e sei.\\
Impianto e reperto & Chi li custodisce, quali parti sono utilizzabili e quali danneggiate, sequestrate o distrutte. Il solo guasto del regolatore non elimina la ricerca.\\
Prove & Chi possiede originali o copie, chi ha ricevuto notizia degli abusi e quali testimonianze restano disponibili.\\
\end{tabularx}

Le combinazioni emergono dalle scelte. Consegnare carte ad Aldren non rende automaticamente complici i PG; preservarle con Miriam non trasferisce il controllo all'Accademia. Distruggere l'impianto non garantisce né esclude che qualcuno riprenda la ricerca altrove.



\textbf{Eventuale contatto:} resta una conclusione da definire separatamente. Non avviene automaticamente durante una lite, un sabotaggio o al termine dell'orologio. La vicenda può concludersi risolvendo la sorte delle persone e della ricerca senza verificare l'ipotesi di Celestine.

\subsection{Collegamenti}
\linkto{Capitolo 20 per combinare gli epiloghi secondo la sorte di Elias e dei pazienti, la custodia dell'impianto e la circolazione delle prove.}

\section{18. Otto scene mobili per riaprire una pista}

Scegli una scena quando ai PG manca un collegamento o una possibilità d'azione. Usala nel luogo indicato, oppure spostala dove sia plausibile incontrare quel PNG. Le scene sono facoltative e non formano una sequenza. Prima di inserirne una, controlla che cosa i PG sanno già e dove si trovano persone e documenti.

Il dettaglio iniziale rende la situazione visibile; parlare, osservare o intervenire permette di approfondirla. Le informazioni ordinarie non richiedono un tiro. Un'opposizione concreta può invece rendere incerto il modo di ottenerle. Se i PG ignorano l'occasione, la scena prosegue secondo gli interessi dei presenti: non trasformarla in un agguato o spostare magicamente gli handout.

\begin{enumerate}[leftmargin=*]
\item \textbf{Il mercante e la bilancia che mente.}

\textit{Quando serve:} manca una strada verso Silas o i PG non sanno con chi discutere il lavoro di Elias.

Al mercato, un venditore di strumenti usati prova a vendere una bilancia. Rimette i pesi in ordine ogni volta che qualcuno tocca il banco, discute con l'ago e cambia argomento prima di finire una frase. Un cliente protesta: i due piatti non tornano alla stessa altezza. Il mercante cerca qualcuno che gli dia ragione.

\textit{Fonte e pista:} aveva lavorato al banco con Elias prima che questi passasse alla Fonderia. Se i PG lo nominano o discutono degli strumenti, ricorda la sua precisione e indica Silas all'Officina delle Lame come vecchio collega ancora in attività. Può spiegare dove trovarlo. Vuole vendere, non custodisce un segreto: ascoltarlo o esaminare la bilancia basta ad avviare la conversazione. L'oggetto è merce ordinaria, non un pezzo del regolatore o una prova dell'esperimento.

\item \textbf{Una cassa, due numeri.}

\textit{Quando serve:} i PG hanno perso il collegamento fra 17 e 17-B, oppure la destinazione del carro.

Vicino al mercato o alla Fonderia, T. Harrow discute con un incaricato dei pagamenti. Il trasporto del giorno 16 risulta sospeso sotto un riferimento e consegnato sotto un altro. Harrow batte il dito sul conteggio: «Una cassa ho caricato. Una ne ho scaricata. Non due viaggi.»

\textit{Fonte e pista:} è il vetturale che ha firmato H05. In privato conferma la ricodifica e la consegna al deposito del Laboratorio Alchemico; indica che alla Fonderia resta la ricevuta. Vuole risolvere il pagamento senza essere accusato di frode. I PG possono ascoltare, chiedere del viaggio o aiutarlo a confrontare i riferimenti. Non possiede informazioni sul tratto sotterraneo o su ciò che fece Elias dopo aver seguito il carro. H05 resta nel luogo già stabilito.

\item \textbf{Marta riconosce una pesatura.}

\textit{Quando serve:} i pesi non sono stati confrontati o i PG non sanno quale lavoratore interpellare.

All'uscita della Fonderia, Marta corregge un operaio che sta copiando una cifra da un foglio: gli fa rileggere il risultato della pesatura invece di usare il valore della riga precedente. Finisce il controllo e raccoglie le sue cose per lasciare il turno.

\textit{Fonte e pista:} ha visto le pesature contestate e conserva H04. Se i PG chiedono di Elias o mostrano i dati della cassa, riconosce il problema e propone di confrontare il foglio originale durante una visita in Fonderia. Il gesto mostra la sua competenza, non introduce una seconda falsificazione. Le confidenze sul turno coperto a Elias restano private e dipendono da come i PG intendono proteggerla.

\item \textbf{La richiesta troppo generica.}

\textit{Quando serve:} l'Archivio sembra inaccessibile o i PG credono di dover cercare fra tutti i registri.

Al banco dell'Archivio, un visitatore chiede tutte le forniture dell'Ospedale. Ilyne gli restituisce la richiesta e gli fa scegliere una consegna, una data o un materiale. Quando l'uomo trova il riferimento sulla ricevuta, lei può finalmente indicare al copista quale registro recuperare.

\textit{Fonte e pista:} Ilyne applica la normale procedura. I PG possono mostrare H05 o H06a, oppure fornire Fonderia, cassa, data e pesi già ricavati da H01/H04. Il confronto apre la consultazione prevista dal nodo. Se cercano Elias, Ilyne segnala anche l'esistenza della disposizione interna H07, mantenendone riservato il contenuto. Il visitatore non ha una commessa segreta da scoprire: serve a rendere comprensibile il funzionamento del banco.

\item \textbf{Il nome che manca.}

\textit{Quando serve:} manca una ragione concreta per parlare con Maud o verificare i trasferimenti.

Nel reparto pubblico dell'Ospedale, Maud chiede di controllare nuovamente un registro. Ricorda il letto e le abitudini del degente, ma chi legge non trova più il suo nome. Accorgendosi che altri possono sentire, l'infermiera interrompe la discussione e torna al lavoro.

\textit{Fonte e pista:} è uno dei due nomi già cancellati fra i sei trasferiti, non una nuova vittima. Se i PG le chiedono di aiutarli a cercare i malati, Maud trova un momento per parlare e mostrare la lista conservata. Può indicare Oren e la porta dell'ala inferiore; non consegna H09 a meno che lui glielo abbia effettivamente affidato. I PG possono attendere, offrirle discrezione o rendere pubblica la contestazione, assumendosi le conseguenze della scelta.

\item \textbf{Nox legge l'ordine fino in fondo.}

\textit{Quando serve:} i PG evitano i Cacciatori perché credono che Elias sia già condannato.

Alla Fonderia o nei pressi del Rifugio, Nox chiede quando Ward sia stato visto. Un interlocutore interpreta l'avviso come un ordine di ucciderlo. Nox torna a leggere H10, soffermandosi sull'intercettazione e sul contenimento, poi precisa che la valutazione spetta alla capitana.

\textit{Fonte e pista:} Nox agisce su incarico di Roland e conosce l'ordine. Può mostrare l'avviso e indicare dove incontrarla per proporre un recupero. La discussione non certifica la lucidità attuale di Elias e non rivela gli esperimenti. Se i PG chiedono dell'aggressione, rimanda alla testimonianza raccolta da Roland. Usa la scena soltanto se Nox è disponibile e quell'incarico è compatibile con quanto è già avvenuto.

\item \textbf{Silas ricorda il contatto.}

\textit{Quando serve:} i PG possiedono la runa o H06b ma non hanno collegato il documento alla visita in Biblioteca.

All'Officina, Silas libera una parte del banco per appoggiare i fogli che i PG stanno consultando. Quando spiegano di cercare il significato del segno, ricorda di avere indirizzato Elias da Miriam e propone di dare loro le indicazioni per raggiungerla.

\textit{Fonte e pista:} Silas procurò il colloquio del giorno 20, ma non ne conosce le conclusioni. Identifica la persona a cui rivolgersi, non traduce la runa. È sufficiente una domanda sul segno o sull'incontro di Elias: non richiede di confessare nuovamente la vendita dei disegni. Se i PG sono già alla Biblioteca, Abel può invece confermare la visita e cercare Miriam secondo la propria scheda. H11 e H11b restano là.

\item \textbf{La cassa che prosegue.}

\textit{Quando serve:} al Laboratorio i PG non sanno dove verificare una consegna o come chiedere dell'ala inferiore.

Durante un'effettiva movimentazione, un lavorante appoggia un collo vicino alle forniture destinate ai tavoli. L'addetto lo ferma, consulta l'inventario sul piano di registrazione e indica il carrello per il trasferimento all'ala inferiore. Poi torna alle altre voci da controllare.

\textit{Fonte e pista:} è una consegna ordinaria, senza un nuovo componente speciale. Il gesto mostra l'uso di H13. Chi chiede del regolatore o della commessa 44-14 può farsi mostrare la riga pertinente; il rimando allo schema di servizio segnala H14 presso l'ascensore. Rimangono validi i controlli della porta. Se non c'è un carico in corso, usa l'addetto che confronta inventario e materiali presenti, senza far partire una spedizione per l'arrivo dei PG.
\end{enumerate}

\section{19. Riferimento rapido al tavolo}

\begin{tabularx}{\linewidth}{>{\bfseries}p{34mm} X}
I PG sono bloccati & Chiedi che cosa stanno cercando di capire. Richiama una prova già ottenuta o scegli una fonte pertinente dalla matrice; usa una scena mobile se serve rendere visibile quella possibilità.\\
Una prova fallisce & Applica una conseguenza all'azione tentata. Se proponi successo a un costo, rendi il costo concreto e lascia scegliere. L'indizio può restare ottenibile da un'altra fonte; non distribuire automaticamente ferite o prove confiscate.\\
Un nodo viene rimandato & Aggiorna orari, persone e controlli secondo gli eventi. I documenti restano dove sono finché qualcuno li sposta; registra autore e destinazione di un eventuale ritiro.\\
Nasce un'alleanza & Concorda aiuto, limiti e impegni. Il PNG rispetta ciò che sa e ciò che può fare; nuove pretese devono avere una ragione emersa in gioco.\\
Inizia un conflitto & Chiarisci obiettivi, posizioni e uscite. Resa, dialogo, protezione dei presenti e fuga restano azioni possibili. Le reazioni dipendono da ciò che ciascuno osserva.\\
Si discute dei testi & Separa fonte, interpretazione e fatto verificato. Miriam può confrontare le prove ricevute; una runa o un beneficio temporaneo non certificano il contatto.\\
Si esplora un ambiente & Descrivi accessi, oggetti, attività e visuali. I dettagli possono essere usati liberamente; richiedi un tiro quando c'è un ostacolo concreto e un esito incerto.\\
Si usano punti Fate & Collega invocazioni e vantaggi ai fatti della scena. Proponi compel pertinenti con una complicazione esplicita; non prepararne uno obbligatorio per ogni PG.\\
Avanza il tempo & Aggiorna l'orologio e le attività con cause leggibili. Precisa il tempo disponibile quando conta, come prima del lavaggio; niente cattura, sabotaggio o contatto automatici.\\
\end{tabularx}

\subsection{Prima della sessione}
\begin{multicols}{2}
\begin{itemize}[label=\tickbox]
\item Segna giorno, ora e stadio dell'orologio.
\item Annota prove e percorsi già scoperti.
\item Controlla chi sa cosa e da quale fonte.
\item Aggiorna posizioni, accordi e incarichi dei PNG.
\item Verifica dove si trovano gli handout.
\item Tieni pronte una o due scene mobili pertinenti.
\item Aggiorna riparo, provviste e intenzioni di Elias.
\item Se si avvicina il finale, colloca i sei pazienti e il personale.
\end{itemize}
\end{multicols}

\section{20. Epiloghi rapidi}

Combina le voci che corrispondono a ciò che è avvenuto: sorte di Elias, assistenza ai sei pazienti, custodia dell'impianto e circolazione delle prove. Gli esiti possono convivere; scegli chi racconta gli eventi in base ai testimoni rimasti.

\begin{description}[leftmargin=5mm,style=nextline]
\item[Elias esce e ritrova Eleanor.]
Se è libero di raggiungerla, possono riunirsi. La malattia rimane; hanno guadagnato tempo e possibilità di scelta.

\item[Elias viene affidato o trattenuto.]
Indica chi ne ha la custodia e quali garanzie sono state ottenute. Gli accordi con Roland, Oren o Aldren contano; la consegna non implica automaticamente esecuzione o guarigione.

\item[Elias muore.]
Testimonianze e documenti determinano ciò che può essere raccontato del suo gesto. Eleanor può conoscere la verità se qualcuno gliela porta; non attribuire alla morte un sabotaggio mai avvenuto.

\item[I pazienti vengono sottratti alle sedute imposte.]
Indica chi assiste tutti e sei e dove. La fine degli abusi non cancella malattia e ricadute; il cambiamento concreto è poter rifiutare ulteriori prove.

\item[Celestine conserva il controllo.]
Se dispone ancora di pazienti, personale e impianto, può tentare di proseguire. Prove diffuse, opposizione interna e vincoli ottenuti dai PG possono limitarla. I benefici restano temporanei.

\item[La ricerca viene sospesa o sequestrata.]
Precisa chi custodisce reperto, macchina e documenti e chi vigila sui pazienti. Una sospensione senza distruzione è un risultato possibile; la posizione futura di Celestine dipende dai fatti resi noti.

\item[Il regolatore o l'impianto vengono distrutti.]
Un guasto al regolatore ferma il lavoro alla portata aumentata; la distruzione più estesa ha conseguenze sulle parti effettivamente colpite. Disegni e risorse conservati possono permettere una ricostruzione, non la rendono certa.

\item[Le prove raggiungono altri studiosi.]
Miriam può favorire contatti e studio dei documenti. La custodia da parte dell'Accademia richiede un accordo effettivo; conservare le prove non autorizza nuove sedute.

\item[Un contatto viene davvero stabilito.]
Usa questa voce soltanto se quel finale è stato definito separatamente. Nessun altro esito lo provoca automaticamente.
\end{description}

\vfill
\begin{center}
\small\itshape
La conclusione prende forma da ciò che i personaggi hanno protetto, cambiato o lasciato nelle mani di altri.
\end{center}

\end{document}
